\section{多变量函数的微分}

\begin{problem}{偏导数在定点的计算}{Partial-Derivative-At-Point}
    求下列各函数在指定点的偏微商：
    \begin{enumerate}
        \item 设 $f(x,y) = x + y - \sqrt{x^2 + y^2}$，求 $f'_x(3,4)$；
        \item 设 $f(x,y) = \sin x^2 y$，求 $f'_x(1, \uppi)$；
        \item 设 $f(x,y) = \ln[xy^2 + yx^2 + \sqrt{1 + (xy^2 + yx^2)^2}]$，求 $f'_x(1, y), f'_y(1, y)$．
    \end{enumerate}
\end{problem}

\begin{solution}
    \ansitem{1}

    对 $x$ 求偏导数得
    \begin{equation}
        f'_x(x, y) = 1 - \frac{x}{\sqrt{x^2 + y^2}}
    \end{equation}
    代入点 $(3,4)$ 可得
    \begin{equation}
        f'_x(3,4) = 1 - \frac{3}{\sqrt{3^2 + 4^2}} = 1 - \frac{3}{5}
    \end{equation}
    \begin{equation}
        \boxed{f'_x(3,4) = \frac{2}{5}}
    \end{equation}

    \ansitem{2}

    对 $x$ 求偏导数得
    \begin{equation}
        f'_x(x, y) = \cos(x^2 y) \cdot 2xy．
    \end{equation}
    将 $x=1, y=\uppi$ 代入上式得
    \begin{equation}
        f'_x(1, \uppi) = 2 \cdot 1 \cdot \uppi \cdot \cos(\uppi) = -2\uppi．
    \end{equation}
    \begin{equation}
        \boxed{f'_x(1, \uppi) = -2\uppi}
    \end{equation}

    \ansitem{3}

    令 $u = xy^2 + yx^2$，则 $f(x,y) = \ln(u + \sqrt{1 + u^2})$．由此可得
    \begin{equation}
        \frac{\partial f}{\partial u} = \frac{1}{\sqrt{1 + u^2}}
    \end{equation}
    分别对 $x, y$ 求偏导数
    \begin{equation}
        \begin{aligned}
            f'_x(x, y) & = \frac{\partial f}{\partial u} \cdot \frac{\partial u}{\partial x} = \frac{y^2 + 2xy}{\sqrt{1 + (xy^2 + yx^2)^2}}, \\
            f'_y(x, y) & = \frac{\partial f}{\partial u} \cdot \frac{\partial u}{\partial y} = \frac{2xy + x^2}{\sqrt{1 + (xy^2 + yx^2)^2}}.
        \end{aligned}
    \end{equation}
    令 $x=1$，则 $u = y^2 + y$，代入得
    \begin{equation}
        \boxed{f'_x(1, y) = \frac{y^2 + 2y}{\sqrt{1 + (y^2 + y)^2}}, \quad f'_y(1, y) = \frac{2y + 1}{\sqrt{1 + (y^2 + y)^2}}}
    \end{equation}
\end{solution}


\begin{problem}{多元函数偏导数计算}{Multivariate-Partial-Derivative}
    求下列各函数对于每个自变量的偏微商：
    \begin{enumerate}
        \item $z = \frac{x \e^y}{y^2}$；
        \item $z = 3^{-\frac{y}{x}}$；
        \item $z = \sin \frac{x}{y} \cos \frac{y}{x}$；
        \item $z = \ln(x + \sqrt{x^2 + y^2})$；
        \item $u = \arctan \frac{x + y}{x - y}$；
        \item $u = \e^{x(x^2 + y^2 + z^2)}$；
        \item $u = x^{y^z}$；
        \item $u = x \e^{-z} + \ln(x + \ln y) + z$．
    \end{enumerate}
\end{problem}

\begin{solution}
    \ansitem{1}

    \begin{equation}
        \frac{\partial z}{\partial x} = \frac{\e^y}{y^2}, \quad \frac{\partial z}{\partial y} = \frac{x \e^y \cdot y^2 - x \e^y \cdot 2y}{y^4} = \frac{x \e^y(y - 2)}{y^3}
    \end{equation}
    \begin{equation}
        \boxed{\frac{\partial z}{\partial x} = \frac{\e^y}{y^2}, \quad \frac{\partial z}{\partial y} = \frac{x(y - 2)\e^y}{y^3}}
    \end{equation}

    \ansitem{2}

    \begin{equation}
        \begin{aligned}
            \frac{\partial z}{\partial x} & = 3^{-\frac{y}{x}} \ln 3 \cdot \frac{\partial}{\partial x}\left( -\frac{y}{x} \right) = \frac{y \ln 3}{x^2} 3^{-\frac{y}{x}}, \\
            \frac{\partial z}{\partial y} & = 3^{-\frac{y}{x}} \ln 3 \cdot \frac{\partial}{\partial y}\left( -\frac{y}{x} \right) = -\frac{\ln 3}{x} 3^{-\frac{y}{x}}.
        \end{aligned}
    \end{equation}
    \begin{equation}
        \boxed{\frac{\partial z}{\partial x} = \frac{y \ln 3}{x^2} 3^{-\frac{y}{x}}, \quad \frac{\partial z}{\partial y} = -\frac{\ln 3}{x} 3^{-\frac{y}{x}}}
    \end{equation}

    \ansitem{3}

    \begin{equation}
        \begin{aligned}
            \frac{\partial z}{\partial x} & = \frac{1}{y} \cos \frac{x}{y} \cos \frac{y}{x} + \sin \frac{x}{y} \left( -\sin \frac{y}{x} \right) \left( -\frac{y}{x^2} \right) = \frac{1}{y} \cos \frac{x}{y} \cos \frac{y}{x} + \frac{y}{x^2} \sin \frac{x}{y} \sin \frac{y}{x},  \\
            \frac{\partial z}{\partial y} & = -\frac{x}{y^2} \cos \frac{x}{y} \cos \frac{y}{x} + \sin \frac{x}{y} \left( -\sin \frac{y}{x} \right) \left( \frac{1}{x} \right) = -\frac{x}{y^2} \cos \frac{x}{y} \cos \frac{y}{x} - \frac{1}{x} \sin \frac{x}{y} \sin \frac{y}{x}.
        \end{aligned}
    \end{equation}
    \begin{equation}
        \boxed{\frac{\partial z}{\partial x} = \frac{1}{y} \cos \frac{x}{y} \cos \frac{y}{x} + \frac{y}{x^2} \sin \frac{x}{y} \sin \frac{y}{x}, \quad \frac{\partial z}{\partial y} = -\frac{x}{y^2} \cos \frac{x}{y} \cos \frac{y}{x} - \frac{1}{x} \sin \frac{x}{y} \sin \frac{y}{x}}
    \end{equation}

    \ansitem{4}

    \begin{equation}
        \begin{aligned}
            \frac{\partial z}{\partial x} & = \frac{1}{x + \sqrt{x^2 + y^2}} \left( 1 + \frac{x}{\sqrt{x^2 + y^2}} \right) = \frac{1}{\sqrt{x^2 + y^2}}, \\
            \frac{\partial z}{\partial y} & = \frac{1}{x + \sqrt{x^2 + y^2}} \cdot \frac{y}{\sqrt{x^2 + y^2}} = \frac{y}{x\sqrt{x^2 + y^2} + x^2 + y^2}.
        \end{aligned}
    \end{equation}
    \begin{equation}
        \boxed{\frac{\partial z}{\partial x} = \frac{1}{\sqrt{x^2 + y^2}}, \quad \frac{\partial z}{\partial y} = \frac{y}{x\sqrt{x^2 + y^2} + x^2 + y^2}}
    \end{equation}

    \ansitem{5}

    \begin{equation}
        \begin{aligned}
            \frac{\partial u}{\partial x} & = \frac{1}{1 + \left( \frac{x+y}{x-y} \right)^2} \cdot \frac{(x-y) - (x+y)}{(x-y)^2} = \frac{-2y}{(x-y)^2 + (x+y)^2} = -\frac{y}{x^2 + y^2},   \\
            \frac{\partial u}{\partial y} & = \frac{1}{1 + \left( \frac{x+y}{x-y} \right)^2} \cdot \frac{(x-y) - (x+y)(-1)}{(x-y)^2} = \frac{2x}{(x-y)^2 + (x+y)^2} = \frac{x}{x^2 + y^2}.
        \end{aligned}
    \end{equation}
    \begin{equation}
        \boxed{\frac{\partial u}{\partial x} = -\frac{y}{x^2 + y^2}, \quad \frac{\partial u}{\partial y} = \frac{x}{x^2 + y^2}}
    \end{equation}

    \ansitem{6}

    \begin{equation}
        \begin{aligned}
            \frac{\partial u}{\partial x} & = \e^{x(x^2 + y^2 + z^2)} \cdot (x^2 + y^2 + z^2 + 2x^2) = (3x^2 + y^2 + z^2) \e^{x(x^2 + y^2 + z^2)}, \\
            \frac{\partial u}{\partial y} & = \e^{x(x^2 + y^2 + z^2)} \cdot (2xy) = 2xy \e^{x(x^2 + y^2 + z^2)},                                   \\
            \frac{\partial u}{\partial z} & = \e^{x(x^2 + y^2 + z^2)} \cdot (2xz) = 2xz \e^{x(x^2 + y^2 + z^2)}.
        \end{aligned}
    \end{equation}
    \begin{equation}
        \boxed{\frac{\partial u}{\partial x} = (3x^2 + y^2 + z^2) \e^{x(x^2 + y^2 + z^2)}, \quad \frac{\partial u}{\partial y} = 2xy \e^{x(x^2 + y^2 + z^2)}, \quad \frac{\partial u}{\partial z} = 2xz \e^{x(x^2 + y^2 + z^2)}}
    \end{equation}

    \ansitem{7}

    利用 $u = \e^{y^z \ln x}$：
    \begin{equation}
        \begin{aligned}
            \frac{\partial u}{\partial x} & = y^z x^{y^z - 1},               \\
            \frac{\partial u}{\partial y} & = x^{y^z} \ln x \cdot z y^{z-1}, \\
            \frac{\partial u}{\partial z} & = x^{y^z} \ln x \cdot y^z \ln y．
        \end{aligned}
    \end{equation}
    \begin{equation}
        \boxed{\frac{\partial u}{\partial x} = y^z x^{y^z - 1}, \quad \frac{\partial u}{\partial y} = z y^{z-1} x^{y^z} \ln x, \quad \frac{\partial u}{\partial z} = y^z x^{y^z} \ln x \cdot \ln y}
    \end{equation}

    \ansitem{8}

    \begin{equation}
        \begin{aligned}
            \frac{\partial u}{\partial x} & = \e^{-z} + \frac{1}{x + \ln y},                                  \\
            \frac{\partial u}{\partial y} & = \frac{1}{x + \ln y} \cdot \frac{1}{y} = \frac{1}{y(x + \ln y)}, \\
            \frac{\partial u}{\partial z} & = -x \e^{-z} + 1．
        \end{aligned}
    \end{equation}
    \begin{equation}
        \boxed{\frac{\partial u}{\partial x} = \e^{-z} + \frac{1}{x + \ln y}, \quad \frac{\partial u}{\partial y} = \frac{1}{y(x + \ln y)}, \quad \frac{\partial u}{\partial z} = 1 - x \e^{-z}}
    \end{equation}
\end{solution}


\begin{problem}{变限积分求导}{Variable-Limit-Integral}
    设 $f(x, y) = \int_1^{x^2 y} \frac{\sin t}{t} \d t$，求 $\frac{\partial f}{\partial x}, \frac{\partial f}{\partial y}$．
\end{problem}

\begin{solution}
    根据多元复合函数求导法则与变限积分求导公式，可得
    \begin{equation}
        \begin{aligned}
            \frac{\partial f}{\partial x} & = \frac{\sin(x^2 y)}{x^2 y} \cdot \frac{\partial}{\partial x}(x^2 y) \\
                                          & = \frac{\sin(x^2 y)}{x^2 y} \cdot 2xy                                \\
                                          & = \frac{2 \sin(x^2 y)}{x},
        \end{aligned}
    \end{equation}
    以及
    \begin{equation}
        \begin{aligned}
            \frac{\partial f}{\partial y} & = \frac{\sin(x^2 y)}{x^2 y} \cdot \frac{\partial}{\partial y}(x^2 y) \\
                                          & = \frac{\sin(x^2 y)}{x^2 y} \cdot x^2                                \\
                                          & = \frac{\sin(x^2 y)}{y}.
        \end{aligned}
    \end{equation}
    故最终结果为
    \begin{equation}
        \boxed{\frac{\partial f}{\partial x} = \frac{2 \sin(x^2 y)}{x}, \quad \frac{\partial f}{\partial y} = \frac{\sin(x^2 y)}{y}}
    \end{equation}
\end{solution}


\begin{problem}{偏导数定义}{Partial-Derivative-Definition}
    设 $f(x, y) = \begin{cases} y \sin \frac{1}{x^2+y^2}, & x^2 + y^2 \neq 0, \\ 0, & x^2 + y^2 = 0 \end{cases}$，考察函数 $f(x, y)$ 在原点 $(0, 0)$ 的偏导数．
\end{problem}

\begin{solution}
    根据偏导数在一点处的定义，计算对 $x$ 的偏导数
    \begin{equation}
        f_x(0, 0) = \lim_{\Delta x \to 0} \frac{f(\Delta x, 0) - f(0, 0)}{\Delta x} = \lim_{\Delta x \to 0} \frac{0 \cdot \sin \frac{1}{\Delta x^2} - 0}{\Delta x} = 0.
    \end{equation}
    计算对 $y$ 的偏导数
    \begin{equation}
        f_y(0, 0) = \lim_{\Delta y \to 0} \frac{f(0, \Delta y) - f(0, 0)}{\Delta y} = \lim_{\Delta y \to 0} \frac{\Delta y \sin \frac{1}{\Delta y^2} - 0}{\Delta y} = \lim_{\Delta y \to 0} \sin \frac{1}{\Delta y^2}.
    \end{equation}
    由于当 $\Delta y \to 0$ 时，$\sin \frac{1}{\Delta y^2}$ 在 $[-1, 1]$ 之间剧烈振荡，极限不存在．
    \begin{equation}
        \boxed{f_x(0, 0) = 0, \quad f_y(0, 0) \text{ 不存在}}
    \end{equation}
\end{solution}


\begin{problem}{连续性与偏导数}{Continuity-Partial-Derivative}
    证明函数 $z = \sqrt{x^2 + y^2}$ 在点 $(0,0)$ 连续但偏导数不存在．
\end{problem}

\begin{myproof}
    考察函数在点 $(0,0)$ 处的极限．由于
    \begin{equation}
        \lim_{(x,y) \to (0,0)} \sqrt{x^2 + y^2} = 0
    \end{equation}
    且 $f(0,0) = 0$，故有 $\lim_{(x,y) \to (0,0)} f(x,y) = f(0,0)$，即函数在点 $(0,0)$ 处连续．
    再考察该点处的偏导数．根据定义有
    \begin{equation}
        \begin{aligned}
            f_x(0,0) & = \lim_{\Delta x \to 0} \frac{f(\Delta x, 0) - f(0,0)}{\Delta x}       \\
                     & = \lim_{\Delta x \to 0} \frac{\sqrt{(\Delta x)^2 + 0^2} - 0}{\Delta x} \\
                     & = \lim_{\Delta x \to 0} \frac{|\Delta x|}{\Delta x}
        \end{aligned}
    \end{equation}
    由于左极限 $\lim_{\Delta x \to 0^-} \frac{|\Delta x|}{\Delta x} = -1$，右极限 $\lim_{\Delta x \to 0^+} \frac{|\Delta x|}{\Delta x} = 1$，两者不相等，故 $f_x(0,0)$ 不存在．
    由对称性可知，偏导数 $f_y(0,0)$ 同样不存在．
\end{myproof}

\begin{problem}{曲线切线与坐标轴夹角}{Tangent-Line-Angle}
    求曲面 $z = \frac{x^2+y^2}{4}$ 与平面 $y = 4$ 的交线在点 $(2, 4, 5)$ 处的切线与 $Ox$ 轴的正向所成的角度．
\end{problem}

\begin{solution}
    曲面与平面的交线方程为
    \begin{equation}
        \begin{cases}
            z = \frac{x^2+y^2}{4} \\
            y = 4
        \end{cases} \implies z = \frac{x^2+16}{4}
    \end{equation}
    将交线看作 $y=4$ 平面内的曲线 $z = z(x)$，该曲线在点 $(2, 4, 5)$ 处切线的斜率为
    \begin{equation}
        k = \frac{\d z}{\d x} \bigg|_{x=2} = \frac{x}{2} \bigg|_{x=2} = 1
    \end{equation}
    切线在空间中的方向向量可取为
    \begin{equation}
        \symbfit{v} = (1, 0, k) = (1, 0, 1)
    \end{equation}
    $Ox$ 轴正向的单位向量为 $\symbfit{i} = (1, 0, 0)$．设切线与 $Ox$ 轴正向所成的角度为 $\alpha$，则有
    \begin{equation}
        \cos \alpha = \frac{\symbfit{v} \cdot \symbfit{i}}{|\symbfit{v}| |\symbfit{i}|} = \frac{1}{\sqrt{1^2 + 0^2 + 1^2} \cdot 1} = \frac{1}{\sqrt{2}}
    \end{equation}
    由于切线角度通常取范围 $[0, \uppi]$，得
    \begin{equation}
        \boxed{\alpha = \frac{\uppi}{4}}
    \end{equation}
\end{solution}


\begin{problem}{空间曲线切线的方向角}{Tangent-Line-Direction}
    求曲线 $\begin{cases} z = \sqrt{x^2 + y^2 + 1}, \\ x = 1 \end{cases}$ 上点 $(1, 1, \sqrt{3})$ 处的切线分别与 $x$ 轴、$y$ 轴、$z$ 轴正向的夹角.
\end{problem}

\begin{solution}
    曲线可由参数 $y$ 表示为参数方程
    \begin{equation}
        \begin{cases}
            x = 1, \\
            y = y, \\
            z = \sqrt{y^2 + 2}.
        \end{cases}
    \end{equation}
    曲线在点 $(1, 1, \sqrt{3})$ 处对应的参数为 $y = 1$. 对各分量求导，得到切向量为
    \begin{equation}
        \symbfit{s} = \left( 0, 1, \frac{y}{\sqrt{y^2 + 2}} \right) \bigg|_{y=1} = \left( 0, 1, \frac{1}{\sqrt{3}} \right).
    \end{equation}
    为计算方向余弦，取其同向的整数倍向量 $\symbfit{s}' = (0, \sqrt{3}, 1)$，其模长为 $|\symbfit{s}'| = \sqrt{0^2 + (\sqrt{3})^2 + 1^2} = 2$.
    设切线与各坐标轴正向的夹角分别为 $\alpha, \beta, \gamma$，则
    \begin{equation}
        \begin{aligned}
            \cos \alpha & = \frac{0}{2} = 0,    \\
            \cos \beta  & = \frac{\sqrt{3}}{2}, \\
            \cos \gamma & = \frac{1}{2}.
        \end{aligned}
    \end{equation}
    由此可得夹角大小
    \begin{equation}
        \boxed{\alpha = \frac{\uppi}{2}, \beta = \frac{\uppi}{6}, \gamma = \frac{\uppi}{3}}
    \end{equation}
\end{solution}

\begin{problem}{热传导方程验证}{Heat-Equation-Verification}
    证明：函数 $u = \frac{1}{\sqrt{t}}\e^{-\frac{x^2}{4t}}$ 满足热传导方程 $\frac{\partial u}{\partial t} = \frac{\partial^2 u}{\partial x^2}$.
\end{problem}

\begin{myproof}
    计算函数 $u$ 对 $t$ 的偏导数
    \begin{equation}
        \begin{aligned}
            \frac{\partial u}{\partial t} & = -\frac{1}{2} t^{-\frac{3}{2}} \e^{-\frac{x^2}{4t}} + t^{-\frac{1}{2}} \e^{-\frac{x^2}{4t}} \cdot \left( \frac{x^2}{4t^2} \right) \\
                                          & = \left( \frac{x^2}{4t^2} - \frac{1}{2t} \right) \frac{1}{\sqrt{t}} \e^{-\frac{x^2}{4t}}                                           \\
                                          & = \left( \frac{x^2 - 2t}{4t^2} \right) u.
        \end{aligned}
    \end{equation}
    接着计算 $u$ 对 $x$ 的一阶偏导数
    \begin{equation}
        \frac{\partial u}{\partial x} = \frac{1}{\sqrt{t}} \e^{-\frac{x^2}{4t}} \cdot \left( -\frac{2x}{4t} \right) = -\frac{x}{2t\sqrt{t}} \e^{-\frac{x^2}{4t}}.
    \end{equation}
    继续对 $x$ 求二阶偏导数
    \begin{equation}
        \begin{aligned}
            \frac{\partial^2 u}{\partial x^2} & = -\frac{1}{2t\sqrt{t}} \e^{-\frac{x^2}{4t}} - \frac{x}{2t\sqrt{t}} \e^{-\frac{x^2}{4t}} \cdot \left( -\frac{2x}{4t} \right) \\
                                              & = -\frac{1}{2t\sqrt{t}} \e^{-\frac{x^2}{4t}} + \frac{x^2}{4t^2\sqrt{t}} \e^{-\frac{x^2}{4t}}                                 \\
                                              & = \left( \frac{x^2 - 2t}{4t^2} \right) \frac{1}{\sqrt{t}} \e^{-\frac{x^2}{4t}}                                               \\
                                              & = \left( \frac{x^2 - 2t}{4t^2} \right) u.
        \end{aligned}
    \end{equation}
    比较两端结果可知
    \begin{equation}
        \boxed{\frac{\partial u}{\partial t} = \frac{\partial^2 u}{\partial x^2}}
    \end{equation}
    原命题得证.
\end{myproof}


\begin{problem}{多元函数高阶偏导数}{Partial-Derivative-Calculation}
    在下列各题中，求 $\frac{\partial^{2} z}{\partial x^{2}}, \frac{\partial^{2} z}{\partial x \partial y}, \frac{\partial^{2} z}{\partial y^{2}}$：
    \begin{enumerate}
        \item $z = \frac{x-y}{x+y}$
        \item $z = \arctan \frac{x+y}{1-xy}$
        \item $z = \ln(x + \sqrt{x^{2} + y^{2}})$
        \item $z = \sin^{2}(ax + by)$
        \item $z = y^{\ln x}$
        \item $z = \arcsin(xy)$
    \end{enumerate}
\end{problem}

\begin{solution}
    \ansitem{1}
    一阶偏导数为
    \begin{equation}
        \frac{\partial z}{\partial x} = \frac{(x+y) - (x-y)}{(x+y)^{2}} = \frac{2y}{(x+y)^{2}}, \quad \frac{\partial z}{\partial y} = \frac{-(x+y) - (x-y)}{(x+y)^{2}} = \frac{-2x}{(x+y)^{2}}.
    \end{equation}
    二阶偏导数为
    \begin{equation}
        \begin{aligned}
            \frac{\partial^{2} z}{\partial x^{2}}        & = \frac{\partial}{\partial x} \left[ 2y(x+y)^{-2} \right] = -4y(x+y)^{-3},                                                                      \\
            \frac{\partial^{2} z}{\partial y^{2}}        & = \frac{\partial}{\partial y} \left[ -2x(x+y)^{-2} \right] = 4x(x+y)^{-3},                                                                      \\
            \frac{\partial^{2} z}{\partial x \partial y} & = \frac{\partial}{\partial y} \left[ \frac{2y}{(x+y)^{2}} \right] = \frac{2(x+y)^{2} - 2y \cdot 2(x+y)}{(x+y)^{4}} = \frac{2x - 2y}{(x+y)^{3}}.
        \end{aligned}
    \end{equation}
    结果为
    \begin{equation}
        \boxed{\frac{\partial^{2} z}{\partial x^{2}} = -\frac{4y}{(x+y)^{3}}, \quad \frac{\partial^{2} z}{\partial x \partial y} = \frac{2(x-y)}{(x+y)^{3}}, \quad \frac{\partial^{2} z}{\partial y^{2}} = \frac{4x}{(x+y)^{3}}.}
    \end{equation}

    \ansitem{2}
    利用恒等式 $\arctan \frac{x+y}{1-xy} = \arctan x + \arctan y + C$（在定义域分片内），得
    \begin{equation}
        \frac{\partial z}{\partial x} = \frac{1}{1+x^{2}}, \quad \frac{\partial z}{\partial y} = \frac{1}{1+y^{2}}.
    \end{equation}
    由此可得
    \begin{equation}
        \boxed{\frac{\partial^{2} z}{\partial x^{2}} = -\frac{2x}{(1+x^{2})^{2}}, \quad \frac{\partial^{2} z}{\partial x \partial y} = 0, \quad \frac{\partial^{2} z}{\partial y^{2}} = -\frac{2y}{(1+y^{2})^{2}}.}
    \end{equation}

    \ansitem{3}
    计算一阶偏导数
    \begin{equation}
        \frac{\partial z}{\partial x} = \frac{1}{x + \sqrt{x^{2}+y^{2}}} \left( 1 + \frac{x}{\sqrt{x^{2}+y^{2}}} \right) = \frac{1}{\sqrt{x^{2}+y^{2}}}.
    \end{equation}
    \begin{equation}
        \frac{\partial z}{\partial y} = \frac{1}{x + \sqrt{x^{2}+y^{2}}} \cdot \frac{y}{\sqrt{x^{2}+y^{2}}}.
    \end{equation}
    计算二阶偏导数
    \begin{equation}
        \begin{aligned}
            \frac{\partial^{2} z}{\partial x^{2}}        & = \frac{\partial}{\partial x} (x^{2}+y^{2})^{-1/2} = -x(x^{2}+y^{2})^{-3/2},                                                                                                                                                            \\
            \frac{\partial^{2} z}{\partial x \partial y} & = \frac{\partial}{\partial y} (x^{2}+y^{2})^{-1/2} = -y(x^{2}+y^{2})^{-3/2},                                                                                                                                                            \\
            \frac{\partial^{2} z}{\partial y^{2}}        & = \frac{\sqrt{x^{2}+y^{2}}(x+\sqrt{x^{2}+y^{2}}) - y \left[ \frac{y}{\sqrt{x^{2}+y^{2}}}(x+\sqrt{x^{2}+y^{2}}) + \sqrt{x^{2}+y^{2}} \frac{y}{\sqrt{x^{2}+y^{2}}} \right]}{(x^{2}+y^{2})(x+\sqrt{x^{2}+y^{2}})^{2}}                      \\
                                                         & = \frac{x\sqrt{x^{2}+y^{2}} + x^{2} + y^{2} - y^{2} - \frac{xy^{2}}{\sqrt{x^{2}+y^{2}}} - y^{2}}{(x^{2}+y^{2})(x+\sqrt{x^{2}+y^{2}})^{2}} = \frac{x^{2}\sqrt{x^{2}+y^{2}} + x^{3} - xy^{2}}{(x^{2}+y^{2})^{3/2}(x+\sqrt{x^{2}+y^{2}})}.
        \end{aligned}
    \end{equation}
    化简后得
    \begin{equation}
        \boxed{\frac{\partial^{2} z}{\partial x^{2}} = -\frac{x}{(x^{2}+y^{2})^{3/2}}, \quad \frac{\partial^{2} z}{\partial x \partial y} = -\frac{y}{(x^{2}+y^{2})^{3/2}}, \quad \frac{\partial^{2} z}{\partial y^{2}} = \frac{x^{3} + x^{2}\sqrt{x^{2}+y^{2}} - xy^{2}}{(x^{2}+y^{2})^{3/2}(x+\sqrt{x^{2}+y^{2}})^{2}}.}
    \end{equation}

    \ansitem{4}
    利用 $\sin^{2} \theta = \frac{1 - \cos 2\theta}{2}$，令 $u = ax+by$，则 $z = \frac{1 - \cos 2u}{2}$．
    \begin{equation}
        \frac{\partial z}{\partial x} = a \sin 2u, \quad \frac{\partial z}{\partial y} = b \sin 2u.
    \end{equation}
    二阶偏导数为
    \begin{equation}
        \boxed{\frac{\partial^{2} z}{\partial x^{2}} = 2a^{2} \cos 2(ax+by), \quad \frac{\partial^{2} z}{\partial x \partial y} = 2ab \cos 2(ax+by), \quad \frac{\partial^{2} z}{\partial y^{2}} = 2b^{2} \cos 2(ax+by).}
    \end{equation}

    \ansitem{5}
    将原式写作 $z = \e^{\ln x \ln y}$．
    \begin{equation}
        \frac{\partial z}{\partial x} = \frac{y^{\ln x} \ln y}{x}, \quad \frac{\partial z}{\partial y} = \ln x \cdot y^{\ln x - 1}.
    \end{equation}
    二阶偏导数为
    \begin{equation}
        \begin{aligned}
            \frac{\partial^{2} z}{\partial x^{2}}        & = \ln y \cdot \frac{\partial}{\partial x} \left( \frac{y^{\ln x}}{x} \right) = \ln y \left[ \frac{\frac{1}{x} y^{\ln x} \ln y \cdot x - y^{\ln x}}{x^{2}} \right] = \frac{y^{\ln x} \ln y (\ln y - 1)}{x^{2}}, \\
            \frac{\partial^{2} z}{\partial y^{2}}        & = \ln x (\ln x - 1) y^{\ln x - 2},                                                                                                                                                                             \\
            \frac{\partial^{2} z}{\partial x \partial y} & = \frac{\partial}{\partial x} (\ln x \cdot y^{\ln x - 1}) = \frac{1}{x} y^{\ln x - 1} + \ln x \cdot y^{\ln x - 1} \ln y \cdot \frac{1}{x} = \frac{y^{\ln x - 1}}{x} (1 + \ln x \ln y).
        \end{aligned}
    \end{equation}
    结果为
    \begin{equation}
        \boxed{\frac{\partial^{2} z}{\partial x^{2}} = \frac{y^{\ln x} \ln y (\ln y - 1)}{x^{2}}, \quad \frac{\partial^{2} z}{\partial x \partial y} = \frac{y^{\ln x - 1}(1 + \ln x \ln y)}{x}, \quad \frac{\partial^{2} z}{\partial y^{2}} = \ln x (\ln x - 1) y^{\ln x - 2}.}
    \end{equation}

    \ansitem{6}
    一阶偏导数为
    \begin{equation}
        \frac{\partial z}{\partial x} = \frac{y}{\sqrt{1-x^{2}y^{2}}}, \quad \frac{\partial z}{\partial y} = \frac{x}{\sqrt{1-x^{2}y^{2}}}.
    \end{equation}
    二阶偏导数为
    \begin{equation}
        \begin{aligned}
            \frac{\partial^{2} z}{\partial x^{2}}        & = y \left[ -\frac{1}{2}(1-x^{2}y^{2})^{-3/2}(-2xy^{2}) \right] = \frac{xy^{3}}{(1-x^{2}y^{2})^{3/2}},                                                                           \\
            \frac{\partial^{2} z}{\partial y^{2}}        & = x \left[ -\frac{1}{2}(1-x^{2}y^{2})^{-3/2}(-2x^{2}y) \right] = \frac{x^{3}y}{(1-x^{2}y^{2})^{3/2}},                                                                           \\
            \frac{\partial^{2} z}{\partial x \partial y} & = \frac{\sqrt{1-x^{2}y^{2}} - y \frac{-2x^{2}y}{2\sqrt{1-x^{2}y^{2}}}}{1-x^{2}y^{2}} = \frac{1-x^{2}y^{2} + x^{2}y^{2}}{(1-x^{2}y^{2})^{3/2}} = \frac{1}{(1-x^{2}y^{2})^{3/2}}.
        \end{aligned}
    \end{equation}
    结果为
    \begin{equation}
        \boxed{\frac{\partial^{2} z}{\partial x^{2}} = \frac{xy^{3}}{(1-x^{2}y^{2})^{3/2}}, \quad \frac{\partial^{2} z}{\partial x \partial y} = \frac{1}{(1-x^{2}y^{2})^{3/2}}, \quad \frac{\partial^{2} z}{\partial y^{2}} = \frac{x^{3}y}{(1-x^{2}y^{2})^{3/2}}.}
    \end{equation}
\end{solution}

\begin{problem}{多元函数高阶偏导数}{Higher-Order-Partial-Derivative}
    设 $u = \e^{xyz}$，求 $\frac{\partial^{3} u}{\partial x \partial y \partial z}, \frac{\partial^{3} u}{\partial x \partial y^{2}}$．
\end{problem}

\begin{solution}
    对 $u = \e^{xyz}$ 依次求导
    \begin{equation}
        \frac{\partial u}{\partial z} = xy \e^{xyz}.
    \end{equation}
    \begin{equation}
        \frac{\partial^{2} u}{\partial y \partial z} = \frac{\partial}{\partial y} (xy \e^{xyz}) = x \e^{xyz} + xy(xz) \e^{xyz} = (x + x^{2}yz) \e^{xyz}.
    \end{equation}
    \begin{equation}
        \frac{\partial^{3} u}{\partial x \partial y \partial z} = \frac{\partial}{\partial x} \left[ (x + x^{2}yz) \e^{xyz} \right] = (1 + 2xyz) \e^{xyz} + (x + x^{2}yz)(yz) \e^{xyz} = (1 + 3xyz + x^{2}y^{2}z^{2}) \e^{xyz}.
    \end{equation}
    另一项计算如下
    \begin{equation}
        \frac{\partial u}{\partial y} = xz \e^{xyz}.
    \end{equation}
    \begin{equation}
        \frac{\partial^{2} u}{\partial y^{2}} = \frac{\partial}{\partial y} (xz \e^{xyz}) = x^{2}z^{2} \e^{xyz}.
    \end{equation}
    \begin{equation}
        \frac{\partial^{3} u}{\partial x \partial y^{2}} = \frac{\partial}{\partial x} (x^{2}z^{2} \e^{xyz}) = 2xz^{2} \e^{xyz} + x^{2}z^{2}(yz) \e^{xyz} = (2xz^{2} + x^{2}yz^{3}) \e^{xyz}.
    \end{equation}
    最终结果为
    \begin{equation}
        \boxed{\frac{\partial^{3} u}{\partial x \partial y \partial z} = (1 + 3xyz + x^{2}y^{2}z^{2}) \e^{xyz}, \quad \frac{\partial^{3} u}{\partial x \partial y^{2}} = (2xz^{2} + x^{2}yz^{3}) \e^{xyz}.}
    \end{equation}
\end{solution}


\begin{problem}{多元函数偏导数}{Partial-Derivatives}
    设 $r = \sqrt{x^2 + y^2 + z^2}$，证明当 $r \neq 0$ 时有
    \begin{enumerate}
        \item $\frac{\partial^2 r}{\partial x^2} + \frac{\partial^2 r}{\partial y^2} + \frac{\partial^2 r}{\partial z^2} = \frac{2}{r};$
        \item $\frac{\partial^2 \ln r}{\partial x^2} + \frac{\partial^2 \ln r}{\partial y^2} + \frac{\partial^2 \ln r}{\partial z^2} = \frac{1}{r^2};$
        \item $\frac{\partial^2}{\partial x^2} \left( \frac{1}{r} \right) + \frac{\partial^2}{\partial y^2} \left( \frac{1}{r} \right) + \frac{\partial^2}{\partial z^2} \left( \frac{1}{r} \right) = 0.$
    \end{enumerate}
\end{problem}

\begin{myproof}
    由于 $r^2 = x^2 + y^2 + z^2$，两边关于 $x$ 求偏导得 $2r \frac{\partial r}{\partial x} = 2x$，故有
    \begin{equation}
        \frac{\partial r}{\partial x} = \frac{x}{r}, \quad \frac{\partial r}{\partial y} = \frac{y}{r}, \quad \frac{\partial r}{\partial z} = \frac{z}{r}. \label{eq:first-order}
    \end{equation}

    \ansitem{1}

    对 \zcref{eq:first-order} 继续求二阶偏导数：
    \begin{equation}
        \begin{aligned}
            \frac{\partial^2 r}{\partial x^2} & = \frac{\partial}{\partial x} \left( \frac{x}{r} \right) = \frac{r - x \cdot \frac{\partial r}{\partial x}}{r^2} \\
                                              & = \frac{r - x \cdot \frac{x}{r}}{r^2} = \frac{r^2 - x^2}{r^3}.
        \end{aligned}
    \end{equation}
    同理可得 $\frac{\partial^2 r}{\partial y^2} = \frac{r^2 - y^2}{r^3}$ 及 $\frac{\partial^2 r}{\partial z^2} = \frac{r^2 - z^2}{r^3}$．将三式相加：
    \begin{equation}
        \begin{aligned}
            \frac{\partial^2 r}{\partial x^2} + \frac{\partial^2 r}{\partial y^2} + \frac{\partial^2 r}{\partial z^2} & = \frac{3r^2 - (x^2 + y^2 + z^2)}{r^3}  \\
                                                                                                                      & = \frac{3r^2 - r^2}{r^3} = \frac{2}{r}.
        \end{aligned}
    \end{equation}
    故原命题成立：
    \begin{equation}
        \boxed{\frac{\partial^2 r}{\partial x^2} + \frac{\partial^2 r}{\partial y^2} + \frac{\partial^2 r}{\partial z^2} = \frac{2}{r}}
    \end{equation}

    \ansitem{2}

    计算 $\ln r$ 的一阶偏导数：
    \begin{equation}
        \frac{\partial \ln r}{\partial x} = \frac{1}{r} \frac{\partial r}{\partial x} = \frac{1}{r} \cdot \frac{x}{r} = \frac{x}{r^2}.
    \end{equation}
    继续求二阶偏导数：
    \begin{equation}
        \begin{aligned}
            \frac{\partial^2 \ln r}{\partial x^2} & = \frac{\partial}{\partial x} \left( \frac{x}{r^2} \right) = \frac{r^2 - x \cdot 2r \frac{\partial r}{\partial x}}{r^4} \\
                                                  & = \frac{r^2 - 2x \cdot \frac{x}{r} \cdot r}{r^4} = \frac{r^2 - 2x^2}{r^4}.
        \end{aligned}
    \end{equation}
    同理可得 $\frac{\partial^2 \ln r}{\partial y^2} = \frac{r^2 - 2y^2}{r^4}$ 及 $\frac{\partial^2 \ln r}{\partial z^2} = \frac{r^2 - 2z^2}{r^4}$．相加得：
    \begin{equation}
        \begin{aligned}
            \sum \frac{\partial^2 \ln r}{\partial x^2} & = \frac{3r^2 - 2(x^2 + y^2 + z^2)}{r^4}    \\
                                                       & = \frac{3r^2 - 2r^2}{r^4} = \frac{1}{r^2}.
        \end{aligned}
    \end{equation}
    故原命题成立：
    \begin{equation}
        \boxed{\frac{\partial^2 \ln r}{\partial x^2} + \frac{\partial^2 \ln r}{\partial y^2} + \frac{\partial^2 \ln r}{\partial z^2} = \frac{1}{r^2}}
    \end{equation}

    \ansitem{3}

    计算 $\frac{1}{r}$ 的一阶偏导数：
    \begin{equation}
        \frac{\partial}{\partial x} \left( \frac{1}{r} \right) = -\frac{1}{r^2} \frac{\partial r}{\partial x} = -\frac{x}{r^3}.
    \end{equation}
    继续求二阶偏导数：
    \begin{equation}
        \begin{aligned}
            \frac{\partial^2}{\partial x^2} \left( \frac{1}{r} \right) & = -\frac{\partial}{\partial x} \left( \frac{x}{r^3} \right) = -\frac{r^3 - x \cdot 3r^2 \frac{\partial r}{\partial x}}{r^6} \\
                                                                       & = -\frac{r^3 - 3x^2 r}{r^6} = \frac{3x^2 - r^2}{r^5}.
        \end{aligned}
    \end{equation}
    利用对称性，将三项相加：
    \begin{equation}
        \begin{aligned}
            \sum \frac{\partial^2}{\partial x^2} \left( \frac{1}{r} \right) & = \frac{3(x^2 + y^2 + z^2) - 3r^2}{r^5} \\
                                                                            & = \frac{3r^2 - 3r^2}{r^5} = 0.
        \end{aligned}
    \end{equation}
    故原命题成立：
    \begin{equation}
        \boxed{\frac{\partial^2}{\partial x^2} \frac{1}{r} + \frac{\partial^2}{\partial y^2} \frac{1}{r} + \frac{\partial^2}{\partial z^2} \frac{1}{r} = 0}
    \end{equation}
\end{myproof}


\begin{problem}{混合偏导数}{Mixed-Partial-Derivatives}
    设
    \begin{equation}
        f(x, y) = \begin{cases} xy \frac{x^2 - y^2}{x^2 + y^2}, & x^2 + y^2 \neq 0, \\ 0, & x^2 + y^2 = 0. \end{cases}
    \end{equation}
    证明：函数的二阶偏导数存在，但所有二阶偏导数（特别是两个混合偏导数）在 $(0, 0)$ 不连续，且 $f_{xy}''(0, 0) \neq f_{yx}''(0, 0)$．
\end{problem}

\begin{myproof}
    当 $(x, y) \neq (0, 0)$ 时，对函数 $f(x, y) = \frac{x^3 y - x y^3}{x^2 + y^2}$ 直接求偏导得
    \begin{equation}
        \label{eq:first-order-partials-1}
        \begin{aligned}
            f_x(x, y) & = \frac{(3x^2 y - y^3)(x^2 + y^2) - 2x(x^3 y - x y^3)}{(x^2 + y^2)^2} = \frac{x^4 y + 4x^2 y^3 - y^5}{(x^2 + y^2)^2}, \\
            f_y(x, y) & = \frac{(x^3 - 3x y^2)(x^2 + y^2) - 2y(x^3 y - x y^3)}{(x^2 + y^2)^2} = \frac{x^5 - 4x^3 y^2 - x y^4}{(x^2 + y^2)^2}.
        \end{aligned}
    \end{equation}
    在原点 $(0, 0)$ 处，利用导数定义计算一阶偏导数
    \begin{equation}
        \label{eq:origin-first-order-2}
        f_x(0, 0) = \lim_{h \to 0} \frac{f(h, 0) - f(0, 0)}{h} = 0, \quad f_y(0, 0) = \lim_{k \to 0} \frac{f(0, k) - f(0, 0)}{k} = 0.
    \end{equation}
    结合 \zcref{eq:first-order-partials-1} 与 \zcref{eq:origin-first-order-2}，计算 $(0, 0)$ 处的二阶偏导数
    \begin{equation}
        \label{eq:second-order-at-origin-3}
        \begin{aligned}
            f_{xx}''(0, 0) & = \lim_{h \to 0} \frac{f_x(h, 0) - f_x(0, 0)}{h} = \lim_{h \to 0} \frac{0 - 0}{h} = 0,   \\
            f_{yy}''(0, 0) & = \lim_{k \to 0} \frac{f_y(0, k) - f_y(0, 0)}{k} = \lim_{k \to 0} \frac{0 - 0}{k} = 0,   \\
            f_{xy}''(0, 0) & = \lim_{k \to 0} \frac{f_x(0, k) - f_x(0, 0)}{k} = \lim_{k \to 0} \frac{-k - 0}{k} = -1, \\
            f_{yx}''(0, 0) & = \lim_{h \to 0} \frac{f_y(h, 0) - f_y(0, 0)}{h} = \lim_{h \to 0} \frac{h - 0}{h} = 1.
        \end{aligned}
    \end{equation}
    由 \zcref{eq:second-order-at-origin-3} 可知二阶偏导数在 $(0, 0)$ 均存在，且满足
    \begin{equation}
        \boxed{f_{xy}''(0, 0) = -1 \neq 1 = f_{yx}''(0, 0)}.
    \end{equation}
    根据二阶混合偏导数相等定理（克莱罗(Clairaut)定理），若混合偏导数在某点连续，则该点处的混合偏导数必与求导次序无关．由上述不等关系推知 $f_{xy}''$ 与 $f_{yx}''$ 在 $(0, 0)$ 处均不连续．同理可验证 $f_{xx}''$ 与 $f_{yy}''$ 在该点的极限与趋近路径相关，故二阶偏导数在 $(0, 0)$ 均不连续．
\end{myproof}


\begin{problem}{全微分计算}{Total-Differential-Calculation}
    求下列函数的微分，或在给定点的微分
    \begin{enumerate}
        \item $z = \ln(x^2 + y^2)$
        \item $z = \frac{xy}{x^2 + y^2}$
        \item $u = \frac{s + t}{s - t}$
        \item $z = \arctan \frac{y}{x}$
        \item $z = \sin(xy)$ 在点 $(0,0)$
        \item $z = x^4 + y^4 - 4x^2y^2$ 在点 $(0,0)$，$(1,1)$
    \end{enumerate}
\end{problem}

\begin{solution}
    \ansitem{1}
    计算偏导数可得 $\frac{\partial z}{\partial x} = \frac{2x}{x^2+y^2}$，$\frac{\partial z}{\partial y} = \frac{2y}{x^2+y^2}$．
    \begin{equation}
        \boxed{\d z = \frac{2x \d x + 2y \d y}{x^2 + y^2}}
    \end{equation}

    \ansitem{2}
    利用商的导数法则计算偏导数：
    \begin{equation}
        \begin{aligned}
            \frac{\partial z}{\partial x} & = \frac{y(x^2 + y^2) - xy(2x)}{(x^2 + y^2)^2} = \frac{y(y^2 - x^2)}{(x^2 + y^2)^2} \\
            \frac{\partial z}{\partial y} & = \frac{x(x^2 + y^2) - xy(2y)}{(x^2 + y^2)^2} = \frac{x(x^2 - y^2)}{(x^2 + y^2)^2}
        \end{aligned}
    \end{equation}
    \begin{equation}
        \boxed{\d z = \frac{y(y^2 - x^2) \d x + x(x^2 - y^2) \d y}{(x^2 + y^2)^2}}
    \end{equation}

    \ansitem{3}
    计算关于 $s$ 和 $t$ 的偏导数：
    \begin{equation}
        \begin{aligned}
            \frac{\partial u}{\partial s} & = \frac{(s - t) - (s + t)}{(s - t)^2} = \frac{-2t}{(s - t)^2}    \\
            \frac{\partial u}{\partial t} & = \frac{(s - t) - (s + t)(-1)}{(s - t)^2} = \frac{2s}{(s - t)^2}
        \end{aligned}
    \end{equation}
    \begin{equation}
        \boxed{\d u = \frac{-2t \d s + 2s \d t}{(s - t)^2}}
    \end{equation}

    \ansitem{4}
    利用复合函数求导法则：
    \begin{equation}
        \d z = \frac{1}{1 + (\frac{y}{x})^2} \d \left( \frac{y}{x} \right) = \frac{x^2}{x^2 + y^2} \frac{x \d y - y \d x}{x^2}
    \end{equation}
    \begin{equation}
        \boxed{\d z = \frac{x \d y - y \d x}{x^2 + y^2}}
    \end{equation}

    \ansitem{5}
    计算在点 $(0,0)$ 处的偏导数：
    \begin{equation}
        \begin{aligned}
            \frac{\partial z}{\partial x}\bigg|_{(0,0)} & = y \cos(xy) \bigg|_{(0,0)} = 0 \\
            \frac{\partial z}{\partial y}\bigg|_{(0,0)} & = x \cos(xy) \bigg|_{(0,0)} = 0
        \end{aligned}
    \end{equation}
    \begin{equation}
        \boxed{\d z|_{(0,0)} = 0}
    \end{equation}

    \ansitem{6}
    函数偏导数为 $\frac{\partial z}{\partial x} = 4x^3 - 8xy^2$，$\frac{\partial z}{\partial y} = 4y^3 - 8x^2y$．
    在点 $(0,0)$ 处，偏导数均为 $0$．
    \begin{equation}
        \boxed{\d z|_{(0,0)} = 0}
    \end{equation}
    在点 $(1,1)$ 处，$\frac{\partial z}{\partial x} = 4 - 8 = -4$，$\frac{\partial z}{\partial y} = 4 - 8 = -4$．
    \begin{equation}
        \boxed{\d z|_{(1,1)} = -4 \d x - 4 \d y}
    \end{equation}
\end{solution}

\begin{problem}{乘积微分法则证明}{Product-Rule-Differential-Proof}
    设 $f$ 和 $g$ 是两个可微函数. 利用微分的定义和定理 9.14, 证明 $\d(fg) = g \d f + f \d g$. (即首先分别写出 $f$ 和 $g$ 的增量的表达式, 分析 $fg$ 的增量的表达式).
\end{problem}

\begin{myproof}
    设 $\symbfit{x} = (x_1, \dots, x_n)$．由于 $f, g$ 在 $\symbfit{x}$ 处可微，其增量可表示为：
    \begin{equation}
        \label{eq:inc-fg}
        \begin{aligned}
            \Delta f & = \d f + \symcal{o}(\rho) \\
            \Delta g & = \d g + \symcal{o}(\rho)
        \end{aligned}
    \end{equation}
    其中 $\rho = \|\Delta \symbfit{x}\|$．分析乘积 $fg$ 的增量：
    \begin{equation}
        \begin{aligned}
            \Delta(fg) & = (f + \Delta f)(g + \Delta g) - fg           \\
                       & = g \Delta f + f \Delta g + \Delta f \Delta g
        \end{aligned}
    \end{equation}
    将 \zcref{eq:inc-fg} 代入上式：
    \begin{equation}
        \begin{aligned}
            \Delta(fg) & = g (\d f + \symcal{o}(\rho)) + f (\d g + \symcal{o}(\rho)) + (\d f + \symcal{o}(\rho))(\d g + \symcal{o}(\rho)) \\
                       & = (g \d f + f \d g) + (g + f) \symcal{o}(\rho) + \d f \d g + \symcal{o}(\rho)^2
        \end{aligned}
    \end{equation}
    由于 $\d f, \d g$ 是关于 $\Delta x_i$ 的线性齐次式，故 $\d f \d g$ 是关于 $\Delta x_i$ 的二次齐次式，满足 $\d f \d g = \symcal{o}(\rho)$．由此得：
    \begin{equation}
        \Delta(fg) = (g \d f + f \d g) + \symcal{o}(\rho)
    \end{equation}
    根据全微分的定义与线性主要部分的唯一性，结论成立．
    \begin{equation}
        \boxed{\d(fg) = g \d f + f \d g}
    \end{equation}
\end{myproof}

\begin{problem}{可微性判定}{Differentiability-Check}
    根据可微的定义，证明：函数 $f(x, y) = \sqrt{|xy|}$ 在原点处不可微.
\end{problem}

\begin{myproof}
    首先计算函数在原点 $(0,0)$ 处的偏导数：
    \begin{equation}
        \begin{aligned}
            f_x(0,0) & = \lim_{\Delta x \to 0} \frac{f(\Delta x, 0) - f(0,0)}{\Delta x} = \lim_{\Delta x \to 0} \frac{0}{ \Delta x} = 0 \\
            f_y(0,0) & = \lim_{\Delta y \to 0} \frac{f(0, \Delta y) - f(0,0)}{\Delta y} = \lim_{\Delta y \to 0} \frac{0}{ \Delta y} = 0
        \end{aligned}
    \end{equation}
    若 $f$ 在 $(0,0)$ 处可微，则必有：
    \begin{equation}
        \lim_{(\Delta x, \Delta y) \to (0,0)} \frac{f(\Delta x, \Delta y) - f(0,0) - [f_x(0,0)\Delta x + f_y(0,0)\Delta y]}{\sqrt{\Delta x^2 + \Delta y^2}} = 0
    \end{equation}
    代入各值考察极限：
    \begin{equation}
        L = \lim_{(\Delta x, \Delta y) \to (0,0)} \frac{\sqrt{|\Delta x \Delta y|}}{\sqrt{\Delta x^2 + \Delta y^2}}
    \end{equation}
    考虑沿直线 $\Delta y = \Delta x \to 0$ 趋近：
    \begin{equation}
        L = \lim_{\Delta x \to 0} \frac{\sqrt{\Delta x^2}}{\sqrt{2\Delta x^2}} = \frac{1}{\sqrt{2}} \neq 0
    \end{equation}
    故该极限不为 $0$，函数在原点处不可微．
\end{myproof}


\begin{problem}{多元函数连续性与可微性}{Multivariable-Continuity-Differentiability}
    证明：函数 $f(x, y) = \begin{cases} \frac{x^2 y}{x^2 + y^2}, & x^2 + y^2 \neq 0, \\ 0, & x^2 + y^2 = 0 \end{cases}$ 在点 $(0, 0)$ 连续且偏导数存在，但在此点不可微．
\end{problem}

\begin{myproof}
    考察 $f(x, y)$ 在 $(0, 0)$ 处的连续性．由不等式 $|xy| \le \frac{1}{2}(x^2 + y^2)$ 可得：
    \begin{equation}
        |f(x, y) - f(0, 0)| = \left| \frac{x^2 y}{x^2 + y^2} \right| = |x| \cdot \frac{|xy|}{x^2 + y^2} \le \frac{1}{2} |x|.
    \end{equation}
    由于 $\lim_{(x, y) \to (0, 0)} \frac{1}{2} |x| = 0$，根据夹逼定理知 $\lim_{(x, y) \to (0, 0)} f(x, y) = 0 = f(0, 0)$，故函数在 $(0, 0)$ 连续．

    利用偏导数定义计算 $(0, 0)$ 处的偏导数：
    \begin{equation}
        \begin{aligned}
            f_x(0, 0) & = \lim_{\Delta x \to 0} \frac{f(\Delta x, 0) - f(0, 0)}{\Delta x} = \lim_{\Delta x \to 0} \frac{0 - 0}{\Delta x} = 0, \\
            f_y(0, 0) & = \lim_{\Delta y \to 0} \frac{f(0, \Delta y) - f(0, 0)}{\Delta y} = \lim_{\Delta y \to 0} \frac{0 - 0}{\Delta y} = 0.
        \end{aligned}
    \end{equation}
    故偏导数 $f_x(0, 0)$ 与 $f_y(0, 0)$ 均存在．

    考查函数在 $(0, 0)$ 处的可微性．令 $\rho = \sqrt{\Delta x^2 + \Delta y^2}$，需判断下式极限是否为 $0$：
    \begin{equation}
        \frac{f(\Delta x, \Delta y) - [f(0, 0) + f_x(0, 0)\Delta x + f_y(0, 0)\Delta y]}{\rho} = \frac{\Delta x^2 \Delta y}{(\Delta x^2 + \Delta y^2)^{3/2}}.
        \label{eq:diff-limit}
    \end{equation}
    取路径 $\Delta y = \Delta x > 0$ 趋于 $0$，则 \zcref{eq:diff-limit} 式变为：
    \begin{equation}
        \lim_{\Delta x \to 0^+} \frac{\Delta x^3}{(2 \Delta x^2)^{3/2}} = \lim_{\Delta x \to 0^+} \frac{\Delta x^3}{2\sqrt{2} \Delta x^3} = \frac{1}{2\sqrt{2}} \neq 0.
    \end{equation}
    故极限不为 $0$，函数在 $(0, 0)$ 不可微．

    \begin{equation}
        \boxed{f(x, y) \text{ 在 } (0, 0) \text{ 连续、偏导数存在但不可微．}}
    \end{equation}
\end{myproof}


\begin{problem}{多元函数连续、可微与偏导数}{Continuity-Differentiability-Partials}
    证明：函数 $f(x,y) = \begin{cases} (x^2 + y^2) \sin \frac{1}{\sqrt{x^2+y^2}}, & x^2 + y^2 \neq 0, \\ 0, & x^2 + y^2 = 0 \end{cases}$ 在点 $(0,0)$ 连续且偏导数存在，但偏导数在点 $(0,0)$ 不连续，而 $f$ 在原点 $(0,0)$ 可微．
\end{problem}

\begin{myproof}
    令 $r = \sqrt{x^2 + y^2}$．由于 $\left| (x^2 + y^2) \sin \frac{1}{\sqrt{x^2+y^2}} \right| \le x^2 + y^2$，则有
    \begin{equation}
        \lim_{(x,y) \to (0,0)} f(x,y) = \lim_{r \to 0} r^2 \sin \frac{1}{r} = 0 = f(0,0).
    \end{equation}
    故 $f(x,y)$ 在点 $(0,0)$ 连续．由偏导数定义：
    \begin{equation}
        \begin{aligned}
            f_x(0,0) & = \lim_{\Delta x \to 0} \frac{f(\Delta x, 0) - f(0,0)}{\Delta x} = \lim_{\Delta x \to 0} \frac{\Delta x^2 \sin \frac{1}{|\Delta x|}}{\Delta x} = 0, \\
            f_y(0,0) & = \lim_{\Delta y \to 0} \frac{f(0, \Delta y) - f(0,0)}{\Delta y} = \lim_{\Delta y \to 0} \frac{\Delta y^2 \sin \frac{1}{|\Delta y|}}{\Delta y} = 0.
        \end{aligned}
    \end{equation}
    故偏导数在 $(0,0)$ 处存在．考察可微性，由于
    \begin{equation}
        \lim_{(\Delta x, \Delta y) \to (0,0)} \frac{f(\Delta x, \Delta y) - f(0,0) - [f_x(0,0) \Delta x + f_y(0,0) \Delta y]}{\sqrt{\Delta x^2 + \Delta y^2}} = \lim_{r \to 0} \frac{r^2 \sin \frac{1}{r}}{r} = 0,
    \end{equation}
    故 $f(x,y)$ 在点 $(0,0)$ 可微．当 $(x,y) \neq (0,0)$ 时，计算 $f_x(x,y)$ 得
    \begin{equation}
        f_x(x,y) = 2x \sin \frac{1}{\sqrt{x^2 + y^2}} - \frac{x}{\sqrt{x^2 + y^2}} \cos \frac{1}{\sqrt{x^2 + y^2}}.
    \end{equation}
    当 $(x,y) \to (0,0)$ 时，第一项趋于 $0$，但第二项极限不存在（震荡），故 $f_x(x,y)$ 在 $(0,0)$ 不连续．同理可知 $f_y(x,y)$ 在 $(0,0)$ 亦不连续．
\end{myproof}

\begin{problem}{多元复合函数求导}{Chain-Rule-Derivatives}
    对下列函数，求出关于 $x$ 和 $y$ 的所有一阶和二阶（含混合）偏导数.
    \begin{enumerate}
        \item $z = u \ln v$，其中 $u = x^2$，$v = \frac{1}{1+y}$；
        \item $z = u \arctan v$，其中 $u = \frac{xy}{x-y}$，$v = x^2y + y - x$.
    \end{enumerate}
\end{problem}

\begin{solution}
    \ansitem{1}

    将中间变量代入，原函数化简为 $z = x^2 \ln \frac{1}{1+y} = -x^2 \ln(1+y)$．
    其一阶偏导数为
    \begin{equation}
        z_x = -2x \ln(1+y), \quad z_y = -\frac{x^2}{1+y}.
    \end{equation}
    对其继续求导，得到二阶偏导数为
    \begin{equation}
        z_{xx} = -2\ln(1+y), \quad z_{yy} = \frac{x^2}{(1+y)^2}, \quad z_{xy} = z_{yx} = -\frac{2x}{1+y}.
    \end{equation}
    汇总最终结果为
    \begin{equation}
        \boxed{
            \begin{aligned}
                z_x    & = -2x \ln(1+y),        \\
                z_y    & = -\frac{x^2}{1+y},    \\
                z_{xx} & = -2\ln(1+y),          \\
                z_{yy} & = \frac{x^2}{(1+y)^2}, \\
                z_{xy} & = -\frac{2x}{1+y}.
            \end{aligned}
        }
    \end{equation}

    \ansitem{2}

    分别计算中间变量 $u, v$ 对 $x, y$ 的一阶与二阶偏导数：
    \begin{equation}
        \begin{aligned}
            u_x    & = -\frac{y^2}{(x-y)^2}, \quad & u_y    & = \frac{x^2}{(x-y)^2},                                           \\
            u_{xx} & = \frac{2y^2}{(x-y)^3}, \quad & u_{yy} & = \frac{2x^2}{(x-y)^3}, \quad & u_{xy} & = -\frac{2xy}{(x-y)^3}, \\
            v_x    & = 2xy - 1, \quad              & v_y    & = x^2 + 1,                                                       \\
            v_{xx} & = 2y, \quad                   & v_{yy} & = 0, \quad                    & v_{xy} & = 2x.
        \end{aligned}
    \end{equation}
    根据多元复合函数求导法则（链式法则与乘积法则），一阶偏导数的结构为
    \begin{equation}
        \begin{aligned}
            z_x & = u_x \arctan v + \frac{u v_x}{1+v^2}, \\
            z_y & = u_y \arctan v + \frac{u v_y}{1+v^2}.
        \end{aligned}
    \end{equation}
    进一步求导得到二阶偏导数的结构式为
    \begin{equation}
        \begin{aligned}
            z_{xx} & = u_{xx} \arctan v + \frac{2 u_x v_x}{1+v^2} + \frac{u v_{xx}}{1+v^2} - \frac{2 u v v_x^2}{(1+v^2)^2},           \\
            z_{yy} & = u_{yy} \arctan v + \frac{2 u_y v_y}{1+v^2} + \frac{u v_{yy}}{1+v^2} - \frac{2 u v v_y^2}{(1+v^2)^2},           \\
            z_{xy} & = u_{xy} \arctan v + \frac{u_x v_y + u_y v_x}{1+v^2} + \frac{u v_{xy}}{1+v^2} - \frac{2 u v v_x v_y}{(1+v^2)^2}.
        \end{aligned}
    \end{equation}
    将所有中间变量及偏导数表达式完全代入上述结构式，最终展开结果如下：
    \begin{equation}
        \boxed{
            \begin{aligned}
                z_x    & = -\frac{y^2 \arctan(x^2y+y-x)}{(x-y)^2} + \frac{xy(2xy-1)}{(x-y)[1+(x^2y+y-x)^2]},                         \\
                z_y    & = \frac{x^2 \arctan(x^2y+y-x)}{(x-y)^2} + \frac{xy(x^2+1)}{(x-y)[1+(x^2y+y-x)^2]},                          \\
                z_{xx} & = \frac{2y^2 \arctan(x^2y+y-x)}{(x-y)^3} - \frac{2y^2(2xy-1)}{(x-y)^2[1+(x^2y+y-x)^2]}                      \\
                       & \quad + \frac{2xy^2}{(x-y)[1+(x^2y+y-x)^2]} - \frac{2xy(x^2y+y-x)(2xy-1)^2}{(x-y)[1+(x^2y+y-x)^2]^2},       \\
                z_{yy} & = \frac{2x^2 \arctan(x^2y+y-x)}{(x-y)^3} + \frac{2x^2(x^2+1)}{(x-y)^2[1+(x^2y+y-x)^2]}                      \\
                       & \quad - \frac{2xy(x^2y+y-x)(x^2+1)^2}{(x-y)[1+(x^2y+y-x)^2]^2},                                             \\
                z_{xy} & = -\frac{2xy \arctan(x^2y+y-x)}{(x-y)^3} + \frac{x^2(2xy-1)-y^2(x^2+1)}{(x-y)^2[1+(x^2y+y-x)^2]}            \\
                       & \quad + \frac{2x^2 y}{(x-y)[1+(x^2y+y-x)^2]} - \frac{2xy(x^2y+y-x)(2xy-1)(x^2+1)}{(x-y)[1+(x^2y+y-x)^2]^2}.
            \end{aligned}
        }
    \end{equation}
\end{solution}


\begin{problem}{复合函数求导}{Composite-Function-Derivative}
    求下列复合函数的偏导数或导数.
    \begin{enumerate}
        \item 设 $u = \e^t + \arctan(t^2 + 1), t = x^y$，求 $\frac{\partial u}{\partial x}, \frac{\partial u}{\partial y}$.
        \item 设 $u = \e^{xyz}, x = rs, y = \frac{r}{s}, z = r^s$，求 $\frac{\partial u}{\partial r}, \frac{\partial u}{\partial s}$.
        \item 设 $u = \ln(x^2 + y^2), x = \e^{t+s+r}, y = 4(s^2 + t^2)$，求 $\frac{\partial u}{\partial r}, \frac{\partial u}{\partial s}, \frac{\partial u}{\partial t}$.
        \item 设 $u = \frac{\e^{ax}(y-z)}{a^2+1}, y = a \sin x, z = \cos x$，求 $\frac{\d u}{\d x}$.
    \end{enumerate}
\end{problem}

\begin{solution}
    \ansitem{1}

    计算中间变量的导数与偏导数：
    \begin{equation}
        \begin{aligned}
            \frac{\d u}{\d t}             & = \e^t + \frac{2t}{1 + (t^2 + 1)^2},                          \\
            \frac{\partial t}{\partial x} & = y x^{y-1}, \quad \frac{\partial t}{\partial y} = x^y \ln x.
        \end{aligned}
    \end{equation}
    由复合函数求导法则，得：
    \begin{equation}
        \begin{aligned}
            \frac{\partial u}{\partial x} & = \frac{\d u}{\d t} \frac{\partial t}{\partial x} = \left[ \e^{x^y} + \frac{2x^y}{1 + (x^{2y} + 1)^2} \right] y x^{y-1}, \\
            \frac{\partial u}{\partial y} & = \frac{\d u}{\d t} \frac{\partial t}{\partial y} = \left[ \e^{x^y} + \frac{2x^y}{1 + (x^{2y} + 1)^2} \right] x^y \ln x.
        \end{aligned}
    \end{equation}
    \begin{equation}
        \boxed{\frac{\partial u}{\partial x} = \left[ \e^{x^y} + \frac{2x^y}{x^{4y} + 2x^{2y} + 2} \right] y x^{y-1}, \quad \frac{\partial u}{\partial y} = \left[ \e^{x^y} + \frac{2x^y}{x^{4y} + 2x^{2y} + 2} \right] x^y \ln x}
    \end{equation}

    \ansitem{2}

    将 $x, y, z$ 代入 $u$ 的表达式中：
    \begin{equation}
        u = \e^{(rs) \cdot \frac{r}{s} \cdot r^s} = \e^{r^2 \cdot r^s} = \e^{r^{s+2}}.
    \end{equation}
    直接求偏导数：
    \begin{equation}
        \begin{aligned}
            \frac{\partial u}{\partial r} & = \e^{r^{s+2}} \cdot (s+2) r^{s+1}, \\
            \frac{\partial u}{\partial s} & = \e^{r^{s+2}} \cdot r^{s+2} \ln r.
        \end{aligned}
    \end{equation}
    \begin{equation}
        \boxed{\frac{\partial u}{\partial r} = (s+2) r^{s+1} \e^{r^{s+2}}, \quad \frac{\partial u}{\partial s} = r^{s+2} \e^{r^{s+2}} \ln r}
    \end{equation}

    \ansitem{3}

    计算各分量的偏导数：
    \begin{equation}
        \begin{aligned}
            \frac{\partial u}{\partial x}                                 & = \frac{2x}{x^2 + y^2}, \quad \frac{\partial u}{\partial y} = \frac{2y}{x^2 + y^2},      \\
            \frac{\partial x}{\partial r} = \frac{\partial x}{\partial s} & = \frac{\partial x}{\partial t} = \e^{t+s+r} = x,                                        \\
            \frac{\partial y}{\partial r}                                 & = 0, \quad \frac{\partial y}{\partial s} = 8s, \quad \frac{\partial y}{\partial t} = 8t.
        \end{aligned}
    \end{equation}
    应用链式法则：
    \begin{equation}
        \begin{aligned}
            \frac{\partial u}{\partial r} & = \frac{\partial u}{\partial x} \frac{\partial x}{\partial r} = \frac{2x^2}{x^2 + y^2},                                                                      \\
            \frac{\partial u}{\partial s} & = \frac{\partial u}{\partial x} \frac{\partial x}{\partial s} + \frac{\partial u}{\partial y} \frac{\partial y}{\partial s} = \frac{2x^2 + 16sy}{x^2 + y^2}, \\
            \frac{\partial u}{\partial t} & = \frac{\partial u}{\partial x} \frac{\partial x}{\partial t} + \frac{\partial u}{\partial y} \frac{\partial y}{\partial t} = \frac{2x^2 + 16ty}{x^2 + y^2}.
        \end{aligned}
    \end{equation}
    \begin{equation}
        \boxed{\frac{\partial u}{\partial r} = \frac{2x^2}{x^2 + y^2}, \quad \frac{\partial u}{\partial s} = \frac{2x^2 + 16sy}{x^2 + y^2}, \quad \frac{\partial u}{\partial t} = \frac{2x^2 + 16ty}{x^2 + y^2}}
    \end{equation}

    \ansitem{4}

    将 $y, z$ 的表达式代入 $u$：
    \begin{equation}
        u = \frac{\e^{ax}(a \sin x - \cos x)}{a^2 + 1}.
    \end{equation}
    对 $x$ 求全导数：
    \begin{equation}
        \begin{aligned}
            \frac{\d u}{\d x} & = \frac{1}{a^2 + 1} \left[ a \e^{ax}(a \sin x - \cos x) + \e^{ax}(a \cos x + \sin x) \right] \\
                              & = \frac{\e^{ax}}{a^2 + 1} \left[ a^2 \sin x - a \cos x + a \cos x + \sin x \right]           \\
                              & = \frac{\e^{ax}}{a^2 + 1} (a^2 + 1) \sin x = \e^{ax} \sin x.
        \end{aligned}
    \end{equation}
    \begin{equation}
        \boxed{\frac{\d u}{\d x} = \e^{ax} \sin x}
    \end{equation}
\end{solution}

\begin{problem}{抽象复合函数偏导数}{Abstract-Composite-Derivative}
    求下列复合函数的偏导数或导数，其中各题中的 $f$ 均有连续的二阶偏导.
    \begin{enumerate}
        \item 设 $u = f(x, y), x = t^3, y = 2t^2$，求 $\frac{\d u}{\d t}$.
        \item 设 $u = f(x, y, z), x = \sin t, y = \cos t, z = \e^t$，求 $\frac{\d u}{\d t}$.
        \item 设 $u = f(x^2 - y^2, \e^{xy})$，求 $\frac{\partial u}{\partial x}, \frac{\partial^2 u}{\partial x \partial y}$.
        \item 设 $u = f(x+y+z, x^2 + y^2 + z^2)$，求 $\frac{\partial u}{\partial x}, \frac{\partial^2 u}{\partial x^2}, \frac{\partial^2 u}{\partial x \partial y}$.
    \end{enumerate}
\end{problem}

\begin{solution}
    \ansitem{1}

    根据复合函数求导法则：
    \begin{equation}
        \frac{\d u}{\d t} = f_1 \frac{\d x}{\d t} + f_2 \frac{\d y}{\d t} = f_1 \cdot 3t^2 + f_2 \cdot 4t.
    \end{equation}
    \begin{equation}
        \boxed{\frac{\d u}{\d t} = 3t^2 f_1 + 4t f_2}
    \end{equation}

    \ansitem{2}

    根据复合函数求导法则：
    \begin{equation}
        \frac{\d u}{\d t} = f_1 \frac{\d x}{\d t} + f_2 \frac{\d y}{\d t} + f_3 \frac{\d z}{\d t} = f_1 \cos t - f_2 \sin t + f_3 \e^t.
    \end{equation}
    \begin{equation}
        \boxed{\frac{\d u}{\d t} = f_1 \cos t - f_2 \sin t + f_3 \e^t}
    \end{equation}

    \ansitem{3}

    设 $u = f(v_1, v_2)$，其中 $v_1 = x^2 - y^2, v_2 = \e^{xy}$.
    \begin{equation}
        \frac{\partial u}{\partial x} = f_1 \cdot 2x + f_2 \cdot y \e^{xy}.
    \end{equation}
    继续对 $y$ 求偏导：
    \begin{equation}
        \begin{aligned}
            \frac{\partial^2 u}{\partial x \partial y} & = \frac{\partial}{\partial y} (2x f_1) + \frac{\partial}{\partial y} (y \e^{xy} f_2)                                                        \\
                                                       & = 2x \left[ f_{11}(-2y) + f_{12}(x \e^{xy}) \right] + (\e^{xy} + xy \e^{xy}) f_2 + y \e^{xy} \left[ f_{21}(-2y) + f_{22}(x \e^{xy}) \right] \\
                                                       & = -4xy f_{11} + 2x^2 \e^{xy} f_{12} + (1 + xy) \e^{xy} f_2 - 2y^2 \e^{xy} f_{21} + xy \e^{2xy} f_{22}.
        \end{aligned}
    \end{equation}
    由于 $f_{12} = f_{21}$：
    \begin{equation}
        \boxed{\frac{\partial^2 u}{\partial x \partial y} = -4xy f_{11} + (2x^2 - 2y^2) \e^{xy} f_{12} + xy \e^{2xy} f_{22} + (1+xy) \e^{xy} f_2}
    \end{equation}

    \ansitem{4}

    设 $u = f(v_1, v_2)$，其中 $v_1 = x+y+z, v_2 = x^2+y^2+z^2$.
    \begin{equation}
        \frac{\partial u}{\partial x} = f_1 \cdot 1 + f_2 \cdot 2x.
    \end{equation}
    \begin{equation}
        \begin{aligned}
            \frac{\partial^2 u}{\partial x^2} & = \frac{\partial}{\partial x} (f_1 + 2x f_2) = (f_{11} + 2x f_{12}) + 2 f_2 + 2x (f_{21} + 2x f_{22}) \\
                                              & = f_{11} + 4x f_{12} + 4x^2 f_{22} + 2f_2.
        \end{aligned}
    \end{equation}
    \begin{equation}
        \begin{aligned}
            \frac{\partial^2 u}{\partial x \partial y} & = \frac{\partial}{\partial y} (f_1 + 2x f_2) = (f_{11} + 2y f_{12}) + 2x (f_{21} + 2y f_{22}) \\
                                                       & = f_{11} + 2(x+y) f_{12} + 4xy f_{22}.
        \end{aligned}
    \end{equation}
    \begin{equation}
        \boxed{\frac{\partial u}{\partial x} = f_1 + 2x f_2, \quad \frac{\partial^2 u}{\partial x^2} = f_{11} + 4x f_{12} + 4x^2 f_{22} + 2f_2, \quad \frac{\partial^2 u}{\partial x \partial y} = f_{11} + 2(x+y) f_{12} + 4xy f_{22}}
    \end{equation}
\end{solution}


\begin{problem}{方向导数}{Directional-Derivative-Basic}
    求函数 $u = xyz$ 在点 $(1, 2, -1)$ 沿方向 $\symbfit{l} = (3, -1, 1)$ 的方向微商．
\end{problem}

\begin{solution}
    计算函数 $u$ 在点 $P(1, 2, -1)$ 处的梯度：
    \begin{equation}
        \nabla u = \left( \frac{\partial u}{\partial x}, \frac{\partial u}{\partial y}, \frac{\partial u}{\partial z} \right) = (yz, xz, xy).
    \end{equation}
    在点 $P(1, 2, -1)$ 处，梯度为：
    \begin{equation}
        \nabla u|_{P} = (-2, -1, 2).
    \end{equation}
    方向 $\symbfit{l} = (3, -1, 1)$ 的单位向量为：
    \begin{equation}
        \symbfit{l}^0 = \frac{\symbfit{l}}{|\symbfit{l}|} = \frac{1}{\sqrt{3^2 + (-1)^2 + 1^2}}(3, -1, 1) = \frac{1}{\sqrt{11}}(3, -1, 1).
    \end{equation}
    则沿该方向的方向微商为：
    \begin{equation}
        \begin{aligned}
            \frac{\partial u}{\partial \symbfit{l}} & = \nabla u|_{P} \cdot \symbfit{l}^0               \\
                                                    & = (-2, -1, 2) \cdot \frac{1}{\sqrt{11}}(3, -1, 1) \\
                                                    & = \frac{1}{\sqrt{11}}(-6 + 1 + 2).
        \end{aligned}
    \end{equation}
    计算得：
    \begin{equation}
        \boxed{\frac{\partial u}{\partial \symbfit{l}} = -\frac{3}{\sqrt{11}}}
    \end{equation}
\end{solution}


\begin{problem}{曲线方向导数}{Curve-Directional-Derivative}
    试求函数 $z = \arctan \frac{y}{x}$ 在圆 $x^2 + y^2 - 2x = 0$ 上一点 $P\left(\frac{1}{2}, \frac{\sqrt{3}}{2}\right)$ 处沿该圆周逆时针方向上的方向微商．
\end{problem}

\begin{solution}
    函数 $z$ 在点 $P$ 处的梯度为：
    \begin{equation}
        \nabla z = \left( \frac{-\frac{y}{x^2}}{1+(\frac{y}{x})^2}, \frac{\frac{1}{x}}{1+(\frac{y}{x})^2} \right) = \left( -\frac{y}{x^2+y^2}, \frac{x}{x^2+y^2} \right).
    \end{equation}
    在点 $P\left(\frac{1}{2}, \frac{\sqrt{3}}{2}\right)$ 处，$x^2+y^2 = 2x = 1$，故：
    \begin{equation}
        \nabla z|_{P} = \left( -\frac{\sqrt{3}}{2}, \frac{1}{2} \right).
    \end{equation}
    圆方程为 $(x-1)^2 + y^2 = 1$，设其参数方程为 $x = 1 + \cos \theta, y = \sin \theta$．在点 $P$ 处有：
    \begin{equation}
        \cos \theta = -\frac{1}{2}, \quad \sin \theta = \frac{\sqrt{3}}{2} \implies \theta = \frac{2\uppi}{3}.
    \end{equation}
    圆周逆时针方向的单位切向量为：
    \begin{equation}
        \symbfit{t} = (-\sin \theta, \cos \theta) = \left( -\frac{\sqrt{3}}{2}, -\frac{1}{2} \right).
    \end{equation}
    则方向微商为：
    \begin{equation}
        \frac{\partial z}{\partial \symbfit{s}} = \nabla z|_{P} \cdot \symbfit{t} = \left( -\frac{\sqrt{3}}{2} \right) \left( -\frac{\sqrt{3}}{2} \right) + \frac{1}{2} \left( -\frac{1}{2} \right) = \frac{3}{4} - \frac{1}{4}.
    \end{equation}
    计算得：
    \begin{equation}
        \boxed{\frac{\partial z}{\partial \symbfit{s}} = \frac{1}{2}}
    \end{equation}
\end{solution}


\begin{problem}{梯度与最大方向导数}{Gradient-Max-Derivative}
    求函数 $u = x^2 + 2y^2 + 3z^2 + xy + 3x - 2y - 6z$ 在点 $(1, 1, -1)$ 的梯度和最大方向微商．
\end{problem}

\begin{solution}
    计算函数 $u$ 的各偏导数：
    \begin{equation}
        \begin{cases}
            u_x = 2x + y + 3, \\
            u_y = 4y + x - 2, \\
            u_z = 6z - 6.
        \end{cases}
    \end{equation}
    在点 $P(1, 1, -1)$ 处：
    \begin{equation}
        \begin{cases}
            u_x|_P = 2(1) + 1 + 3 = 6, \\
            u_y|_P = 4(1) + 1 - 2 = 3, \\
            u_z|_P = 6(-1) - 6 = -12.
        \end{cases}
    \end{equation}
    故梯度为：
    \begin{equation}
        \boxed{\nabla u|_P = (6, 3, -12)}
    \end{equation}
    函数在点 $P$ 处沿梯度方向的方向微商取得最大值，即：
    \begin{equation}
        \left( \frac{\partial u}{\partial \symbfit{l}} \right)_{\max} = |\nabla u|_P = \sqrt{6^2 + 3^2 + (-12)^2} = \sqrt{189}.
    \end{equation}
    化简得：
    \begin{equation}
        \boxed{\left( \frac{\partial u}{\partial \symbfit{l}} \right)_{\max} = 3\sqrt{21}}
    \end{equation}
\end{solution}


\begin{problem}{梯度计算}{Gradient-Calculation}
    设 $\symbfit{r} = x\symbfit{i} + y\symbfit{j} + z\symbfit{k}$，$r = |\symbfit{r}|$，试求
    \begin{enumerate}
        \item $\operatorname{grad} \frac{1}{r^2}$；
        \item $\operatorname{grad} \ln r$．
    \end{enumerate}
\end{problem}

\begin{solution}
    \ansitem{1}

    由于 $\operatorname{grad} r = \frac{\symbfit{r}}{r}$，根据复合函数求导法则有
    \begin{equation}
        \begin{aligned}
            \operatorname{grad} \frac{1}{r^2} & = -\frac{2}{r^3} \operatorname{grad} r       \\
                                              & = -\frac{2}{r^3} \cdot \frac{\symbfit{r}}{r} \\
                                              & = -\frac{2 \symbfit{r}}{r^4}.
        \end{aligned}
    \end{equation}
    \begin{equation}
        \boxed{\operatorname{grad} \frac{1}{r^2} = -\frac{2 \symbfit{r}}{r^4}}
    \end{equation}

    \ansitem{2}

    同理可得
    \begin{equation}
        \begin{aligned}
            \operatorname{grad} \ln r & = \frac{1}{r} \operatorname{grad} r       \\
                                      & = \frac{1}{r} \cdot \frac{\symbfit{r}}{r} \\
                                      & = \frac{\symbfit{r}}{r^2}.
        \end{aligned}
    \end{equation}
    \begin{equation}
        \boxed{\operatorname{grad} \ln r = \frac{\symbfit{r}}{r^2}}
    \end{equation}
\end{solution}

\begin{problem}{多元复合函数偏导数}{Partial-Derivative-Chain-Rule}
    设 $u = f(t)$，$t = \varphi(xy, x + y)$，其中 $f, \varphi$ 分别具有连续的二阶导数及偏导数，求 $\frac{\partial u}{\partial x}, \frac{\partial u}{\partial y}, \frac{\partial^2 u}{\partial x \partial y}$．
\end{problem}

\begin{solution}
    记 $v = xy$，$w = x + y$，则 $t = \varphi(v, w)$．由链式法则得
    \begin{equation}
        \frac{\partial u}{\partial x} = \frac{\d f}{\d t} \left( \frac{\partial \varphi}{\partial v} \frac{\partial v}{\partial x} + \frac{\partial \varphi}{\partial w} \frac{\partial w}{\partial x} \right) = f'(t) (y \varphi_1 + \varphi_2)
    \end{equation}
    \begin{equation}
        \frac{\partial u}{\partial y} = \frac{\d f}{\d t} \left( \frac{\partial \varphi}{\partial v} \frac{\partial v}{\partial y} + \frac{\partial \varphi}{\partial w} \frac{\partial w}{\partial y} \right) = f'(t) (x \varphi_1 + \varphi_2)
    \end{equation}
    对 $\frac{\partial u}{\partial y}$ 关于 $x$ 求偏导
    \begin{equation}
        \begin{aligned}
            \frac{\partial^2 u}{\partial x \partial y} & = \frac{\partial}{\partial x} [f'(t) (x \varphi_1 + \varphi_2)]                                                                                                            \\
                                                       & = \frac{\partial f'(t)}{\partial x} (x \varphi_1 + \varphi_2) + f'(t) \frac{\partial}{\partial x} (x \varphi_1 + \varphi_2)                                                \\
                                                       & = f''(t) \frac{\partial t}{\partial x} (x \varphi_1 + \varphi_2) + f'(t) [\varphi_1 + x (\varphi_{11} y + \varphi_{12} \cdot 1) + (\varphi_{21} y + \varphi_{22} \cdot 1)] \\
                                                       & = f''(t) (y \varphi_1 + \varphi_2) (x \varphi_1 + \varphi_2) + f'(t) (\varphi_1 + xy \varphi_{11} + (x + y) \varphi_{12} + \varphi_{22}).
        \end{aligned}
    \end{equation}
    \begin{equation}
        \boxed{\frac{\partial^2 u}{\partial x \partial y} = f''(t) (y \varphi_1 + \varphi_2) (x \varphi_1 + \varphi_2) + f'(t) (\varphi_1 + xy \varphi_{11} + (x + y) \varphi_{12} + \varphi_{22})}
    \end{equation}
\end{solution}

\begin{problem}{偏导数恒等式证明}{Partial-Derivative-Identity}
    设 $z = f(xy)$，$f$ 为可微函数．证明 $x \frac{\partial z}{\partial x} - y \frac{\partial z}{\partial y} = 0$．
\end{problem}

\begin{myproof}
    令 $u = xy$，则 $z = f(u)$．计算其偏导数
    \begin{equation}
        \frac{\partial z}{\partial x} = f'(u) \frac{\partial u}{\partial x} = y f'(xy)
    \end{equation}
    \begin{equation}
        \frac{\partial z}{\partial y} = f'(u) \frac{\partial u}{\partial y} = x f'(xy)
    \end{equation}
    代入原式左端
    \begin{equation}
        \begin{aligned}
            x \frac{\partial z}{\partial x} - y \frac{\partial z}{\partial y} & = x (y f'(xy)) - y (x f'(xy)) \\
                                                                              & = xy f'(xy) - xy f'(xy)       \\
                                                                              & = 0.
        \end{aligned}
    \end{equation}
    得证．
\end{myproof}


\begin{problem}{复合函数求导}{Composite-Function-Differentiation}
    设 $z = f\left(\ln x + \frac{1}{y}\right)$，$f$ 为可微函数．证明 $x \frac{\partial z}{\partial x} + y^2 \frac{\partial z}{\partial y} = 0$．
\end{problem}

\begin{myproof}
    令 $u = \ln x + \frac{1}{y}$，根据复合函数求导法则，分别计算 $z$ 对 $x$ 和 $y$ 的偏导数：
    \begin{equation} \label{eq:z-partial}
        \begin{aligned}
            \frac{\partial z}{\partial x} & = f'(u) \cdot \frac{\partial u}{\partial x} = f'(u) \cdot \frac{1}{x},                   \\
            \frac{\partial z}{\partial y} & = f'(u) \cdot \frac{\partial u}{\partial y} = f'(u) \cdot \left( -\frac{1}{y^2} \right).
        \end{aligned}
    \end{equation}
    将 \zcref{eq:z-partial} 中的结果代入待证等式的左端：
    \begin{equation}
        x \frac{\partial z}{\partial x} + y^2 \frac{\partial z}{\partial y} = x \cdot \left( \frac{1}{x} f'(u) \right) + y^2 \cdot \left( -\frac{1}{y^2} f'(u) \right) = f'(u) - f'(u) = 0.
    \end{equation}
    由此可知，原等式成立．
    \begin{equation}
        \boxed{x \frac{\partial z}{\partial x} + y^2 \frac{\partial z}{\partial y} = 0}
    \end{equation}
\end{myproof}


\begin{problem}{波动方程}{Wave-Equation-Verification}
    证明函数 $u = \varphi(x - at) + \psi(x + at)$ 满足波动方程
    \begin{equation}
        \frac{\partial^2 u}{\partial t^2} = a^2 \frac{\partial^2 u}{\partial x^2},
    \end{equation}
    其中 $\varphi, \psi$ 有连续的二阶微商．
\end{problem}

\begin{myproof}
    令 $\xi = x - at$，$\eta = x + at$．对 $u$ 关于 $x$ 求一阶及二阶偏导数：
    \begin{equation} \label{eq:u-xx}
        \begin{aligned}
            \frac{\partial u}{\partial x}     & = \varphi'(\xi) \frac{\partial \xi}{\partial x} + \psi'(\eta) \frac{\partial \eta}{\partial x} = \varphi'(\xi) + \psi'(\eta),     \\
            \frac{\partial^2 u}{\partial x^2} & = \varphi''(\xi) \frac{\partial \xi}{\partial x} + \psi''(\eta) \frac{\partial \eta}{\partial x} = \varphi''(\xi) + \psi''(\eta).
        \end{aligned}
    \end{equation}
    对 $u$ 关于 $t$ 求一阶及二阶偏导数：
    \begin{equation} \label{eq:u-tt}
        \begin{aligned}
            \frac{\partial u}{\partial t}     & = \varphi'(\xi) \frac{\partial \xi}{\partial t} + \psi'(\eta) \frac{\partial \eta}{\partial t} = -a \varphi'(\xi) + a \psi'(\eta),                                                                   \\
            \frac{\partial^2 u}{\partial t^2} & = -a \varphi''(\xi) \frac{\partial \xi}{\partial t} + a \psi''(\eta) \frac{\partial \eta}{\partial t} = (-a)^2 \varphi''(\xi) + a^2 \psi''(\eta) = a^2 \left[ \varphi''(\xi) + \psi''(\eta) \right].
        \end{aligned}
    \end{equation}
    比较 \zcref{eq:u-xx} 与 \zcref{eq:u-tt} 可得：
    \begin{equation}
        \frac{\partial^2 u}{\partial t^2} = a^2 \left[ \varphi''(\xi) + \psi''(\eta) \right] = a^2 \frac{\partial^2 u}{\partial x^2}.
    \end{equation}
    结论得证．
    \begin{equation}
        \boxed{\frac{\partial^2 u}{\partial t^2} = a^2 \frac{\partial^2 u}{\partial x^2}}
    \end{equation}
\end{myproof}


\begin{problem}{极坐标变换}{Polar-Coordinate-Substitution}
    若 $u = F(x, y)$，$F$ 任意二阶偏导存在，而 $x = r \cos \varphi$，$y = r \sin \varphi$．证明：
    \begin{equation}
        \left(\frac{\partial u}{\partial r}\right)^2 + \left(\frac{1}{r} \frac{\partial u}{\partial \varphi}\right)^2 = \left(\frac{\partial u}{\partial x}\right)^2 + \left(\frac{\partial u}{\partial y}\right)^2.
    \end{equation}
\end{problem}

\begin{myproof}
    根据多元复合函数求导的链式法则，有
    \begin{equation}
        \label{eq:chain-rule-1}
        \begin{aligned}
            \frac{\partial u}{\partial r}       & = \frac{\partial u}{\partial x} \frac{\partial x}{\partial r} + \frac{\partial u}{\partial y} \frac{\partial y}{\partial r} = \frac{\partial u}{\partial x} \cos \varphi + \frac{\partial u}{\partial y} \sin \varphi,                  \\
            \frac{\partial u}{\partial \varphi} & = \frac{\partial u}{\partial x} \frac{\partial x}{\partial \varphi} + \frac{\partial u}{\partial y} \frac{\partial y}{\partial \varphi} = -\frac{\partial u}{\partial x} r \sin \varphi + \frac{\partial u}{\partial y} r \cos \varphi.
        \end{aligned}
    \end{equation}
    对 $\frac{\partial u}{\partial \varphi}$ 式两边同除以 $r$，得到
    \begin{equation}
        \label{eq:chain-rule-2}
        \frac{1}{r} \frac{\partial u}{\partial \varphi} = -\frac{\partial u}{\partial x} \sin \varphi + \frac{\partial u}{\partial y} \cos \varphi.
    \end{equation}
    将 \zcref{eq:chain-rule-1} 与 \zcref{eq:chain-rule-2} 的等号两边分别平方并相加，得
    \begin{equation}
        \begin{aligned}
            \left(\frac{\partial u}{\partial r}\right)^2 + \left(\frac{1}{r} \frac{\partial u}{\partial \varphi}\right)^2 & = \left( \frac{\partial u}{\partial x} \cos \varphi + \frac{\partial u}{\partial y} \sin \varphi \right)^2 + \left( -\frac{\partial u}{\partial x} \sin \varphi + \frac{\partial u}{\partial y} \cos \varphi \right)^2 \\
                                                                                                                          & = \left( \frac{\partial u}{\partial x} \right)^2 \left( \cos^2 \varphi + \sin^2 \varphi \right) + \left( \frac{\partial u}{\partial y} \right)^2 \left( \sin^2 \varphi + \cos^2 \varphi \right)                        \\
                                                                                                                          & = \left( \frac{\partial u}{\partial x} \right)^2 + \left( \frac{\partial u}{\partial y} \right)^2.
        \end{aligned}
    \end{equation}
    由此证得
    \begin{equation}
        \boxed{\left(\frac{\partial u}{\partial r}\right)^2 + \left(\frac{1}{r} \frac{\partial u}{\partial \varphi}\right)^2 = \left(\frac{\partial u}{\partial x}\right)^2 + \left(\frac{\partial u}{\partial y}\right)^2}
    \end{equation}
\end{myproof}

\begin{problem}{偏微分方程变量代换}{PDE-Variable-Substitution}
    试证：方程 $\frac{\partial^2 u}{\partial x^2} + 2 \frac{\partial^2 u}{\partial x \partial y} - 3 \frac{\partial^2 u}{\partial y^2} + 2 \frac{\partial u}{\partial x} + 6 \frac{\partial u}{\partial y} = 0$ 经变化 $\xi = x + y, \eta = 3x - y$ 后变成 $\frac{\partial^2 u}{\partial \eta \partial \xi} + \frac{1}{2} \frac{\partial u}{\partial \xi} = 0$．（其中二阶偏导数均连续）
\end{problem}

\begin{myproof}
    根据变量代换关系，计算一阶偏导数算子：
    \begin{equation}
        \label{eq:operator-1}
        \begin{aligned}
            \frac{\partial}{\partial x} & = \frac{\partial \xi}{\partial x} \frac{\partial}{\partial \xi} + \frac{\partial \eta}{\partial x} \frac{\partial}{\partial \eta} = \frac{\partial}{\partial \xi} + 3 \frac{\partial}{\partial \eta}, \\
            \frac{\partial}{\partial y} & = \frac{\partial \xi}{\partial y} \frac{\partial}{\partial \xi} + \frac{\partial \eta}{\partial y} \frac{\partial}{\partial \eta} = \frac{\partial}{\partial \xi} - \frac{\partial}{\partial \eta}.
        \end{aligned}
    \end{equation}
    计算二阶偏导数算子：
    \begin{equation}
        \begin{aligned}
            \frac{\partial^2}{\partial x^2}          & = \left( \frac{\partial}{\partial \xi} + 3 \frac{\partial}{\partial \eta} \right)^2 = \frac{\partial^2}{\partial \xi^2} + 6 \frac{\partial^2}{\partial \xi \partial \eta} + 9 \frac{\partial^2}{\partial \eta^2},                                                                             \\
            \frac{\partial^2}{\partial y^2}          & = \left( \frac{\partial}{\partial \xi} - \frac{\partial}{\partial \eta} \right)^2 = \frac{\partial^2}{\partial \xi^2} - 2 \frac{\partial^2}{\partial \xi \partial \eta} + \frac{\partial^2}{\partial \eta^2},                                                                                 \\
            \frac{\partial^2}{\partial x \partial y} & = \left( \frac{\partial}{\partial \xi} + 3 \frac{\partial}{\partial \eta} \right) \left( \frac{\partial}{\partial \xi} - \frac{\partial}{\partial \eta} \right) = \frac{\partial^2}{\partial \xi^2} + 2 \frac{\partial^2}{\partial \xi \partial \eta} - 3 \frac{\partial^2}{\partial \eta^2}.
        \end{aligned}
    \end{equation}
    将上述算子代入原方程：
    \begin{equation}
        \begin{aligned}
             & \left( \frac{\partial^2 u}{\partial \xi^2} + 6 \frac{\partial^2 u}{\partial \xi \partial \eta} + 9 \frac{\partial^2 u}{\partial \eta^2} \right) + 2 \left( \frac{\partial^2 u}{\partial \xi^2} + 2 \frac{\partial^2 u}{\partial \xi \partial \eta} - 3 \frac{\partial^2 u}{\partial \eta^2} \right) - 3 \left( \frac{\partial^2 u}{\partial \xi^2} - 2 \frac{\partial^2 u}{\partial \xi \partial \eta} + \frac{\partial^2 u}{\partial \eta^2} \right) \\
             & + 2 \left( \frac{\partial u}{\partial \xi} + 3 \frac{\partial u}{\partial \eta} \right) + 6 \left( \frac{\partial u}{\partial \xi} - \frac{\partial u}{\partial \eta} \right) = 0.
        \end{aligned}
    \end{equation}
    化简各项系数：
    \begin{equation}
        \begin{aligned}
            \frac{\partial^2 u}{\partial \xi^2}             & : 1 + 2 - 3 = 0,  \\
            \frac{\partial^2 u}{\partial \eta^2}            & : 9 - 6 - 3 = 0,  \\
            \frac{\partial^2 u}{\partial \xi \partial \eta} & : 6 + 4 + 6 = 16, \\
            \frac{\partial u}{\partial \xi}                 & : 2 + 6 = 8,      \\
            \frac{\partial u}{\partial \eta}                & : 6 - 6 = 0.
        \end{aligned}
    \end{equation}
    故方程变为 $16 \frac{\partial^2 u}{\partial \xi \partial \eta} + 8 \frac{\partial u}{\partial \xi} = 0$，即
    \begin{equation}
        \boxed{\frac{\partial^2 u}{\partial \eta \partial \xi} + \frac{1}{2} \frac{\partial u}{\partial \xi} = 0}
    \end{equation}
\end{myproof}


\begin{problem}{偏导数坐标变换}{PDE-Transformation}
    试证：方程 $\frac{\partial^2 u}{\partial x^2} + 2 \frac{\partial^2 u}{\partial x \partial y} \cos x - \frac{\partial^2 u}{\partial y^2} \sin^2 x - \frac{\partial u}{\partial y} \sin x = 0$ 经变换 $\xi = x - \sin x + y, \eta = x + \sin x - y$ 后变成 $\frac{\partial^2 u}{\partial \xi \partial \eta} = 0$．（其中二阶偏导数均连续）
\end{problem}

\begin{myproof}
    根据变换关系，计算一阶偏导数算子：
    \begin{equation}
        \begin{aligned}
            \frac{\partial}{\partial x} & = \frac{\partial \xi}{\partial x} \frac{\partial}{\partial \xi} + \frac{\partial \eta}{\partial x} \frac{\partial}{\partial \eta} = (1 - \cos x) \frac{\partial}{\partial \xi} + (1 + \cos x) \frac{\partial}{\partial \eta}, \\
            \frac{\partial}{\partial y} & = \frac{\partial \xi}{\partial y} \frac{\partial}{\partial \xi} + \frac{\partial \eta}{\partial y} \frac{\partial}{\partial \eta} = \frac{\partial}{\partial \xi} - \frac{\partial}{\partial \eta}.
        \end{aligned}
    \end{equation}
    由此可得 $u_y = u_\xi - u_\eta$，进一步计算二阶偏导数：
    \begin{equation}
        \begin{aligned}
            \frac{\partial^2 u}{\partial x^2}          & = \frac{\partial}{\partial x} \left[ (1 - \cos x) u_\xi + (1 + \cos x) u_\eta \right]                                                                      \\
                                                       & = \sin x (u_\xi - u_\eta) + (1 - \cos x) \frac{\partial u_\xi}{\partial x} + (1 + \cos x) \frac{\partial u_\eta}{\partial x}                               \\
                                                       & = \sin x (u_\xi - u_\eta) + (1 - \cos x)^2 u_{\xi\xi} + 2(1 - \cos^2 x) u_{\xi\eta} + (1 + \cos x)^2 u_{\eta\eta},                                         \\
            \frac{\partial^2 u}{\partial x \partial y} & = \frac{\partial}{\partial x} (u_\xi - u_\eta) = (1 - \cos x) u_{\xi\xi} + (1 + \cos x) u_{\xi\eta} - (1 - \cos x) u_{\eta\xi} - (1 + \cos x) u_{\eta\eta} \\
                                                       & = (1 - \cos x) u_{\xi\xi} + 2\cos x u_{\xi\eta} - (1 + \cos x) u_{\eta\eta},                                                                               \\
            \frac{\partial^2 u}{\partial y^2}          & = \frac{\partial}{\partial y} (u_\xi - u_\eta) = u_{\xi\xi} - 2u_{\xi\eta} + u_{\eta\eta}.
        \end{aligned}
    \end{equation}
    将上述结果代入原方程左端：
    \begin{equation}
        \begin{aligned}
            L & = \left[ \sin x (u_\xi - u_\eta) + (1 - \cos x)^2 u_{\xi\xi} + 2\sin^2 x u_{\xi\eta} + (1 + \cos x)^2 u_{\eta\eta} \right] \\
              & \quad + 2\cos x \left[ (1 - \cos x) u_{\xi\xi} + 2\cos x u_{\xi\eta} - (1 + \cos x) u_{\eta\eta} \right]                   \\
              & \quad - \sin^2 x \left[ u_{\xi\xi} - 2u_{\xi\eta} + u_{\eta\eta} \right] - \sin x (u_\xi - u_\eta)                         \\
              & = u_{\xi\xi} \left[ (1 - \cos x)^2 + 2\cos x (1 - \cos x) - \sin^2 x \right]                                               \\
              & \quad + u_{\eta\eta} \left[ (1 + \cos x)^2 - 2\cos x (1 + \cos x) - \sin^2 x \right]                                       \\
              & \quad + u_{\xi\eta} \left[ 2\sin^2 x + 4\cos^2 x + 2\sin^2 x \right].
        \end{aligned}
    \end{equation}
    化简各项系数：
    \begin{equation}
        \begin{aligned}
            u_{\xi\xi}: \quad   & 1 - 2\cos x + \cos^2 x + 2\cos x - 2\cos^2 x - \sin^2 x = 1 - \cos^2 x - \sin^2 x = 0, \\
            u_{\eta\eta}: \quad & 1 + 2\cos x + \cos^2 x - 2\cos x - 2\cos^2 x - \sin^2 x = 1 - \cos^2 x - \sin^2 x = 0, \\
            u_{\xi\eta}: \quad  & 4\sin^2 x + 4\cos^2 x = 4.
        \end{aligned}
    \end{equation}
    因此方程化为 $4 u_{\xi\eta} = 0$，即：
    \begin{equation}
        \frac{\partial^2 u}{\partial \xi \partial \eta} = 0.
    \end{equation}
\end{myproof}


\begin{problem}{偏导数变量替换}{Variable-Substitution}
    设变换 $\begin{cases} u = x - 2y, \\ v = x + ay \end{cases}$ 可把方程 $6 \frac{\partial^2 z}{\partial x^2} + \frac{\partial^2 z}{\partial x \partial y} - \frac{\partial^2 z}{\partial y^2} = 0$ 简化为 $\frac{\partial^2 z}{\partial u \partial v} = 0$．求常数 $a$．（其中二阶偏导数均连续）
\end{problem}

\begin{solution}
    计算一阶与二阶偏导数算子的对应关系：
    \begin{equation}
        \frac{\partial}{\partial x} = \frac{\partial}{\partial u} + \frac{\partial}{\partial v}, \quad \frac{\partial}{\partial y} = -2\frac{\partial}{\partial u} + a\frac{\partial}{\partial v}.
    \end{equation}
    代入计算二阶偏导数：
    \begin{equation}
        \begin{aligned}
            \frac{\partial^2 z}{\partial x^2}          & = z_{uu} + 2z_{uv} + z_{vv},                                                                                                                 \\
            \frac{\partial^2 z}{\partial x \partial y} & = \left( \frac{\partial}{\partial u} + \frac{\partial}{\partial v} \right) \left( -2z_u + az_v \right) = -2z_{uu} + (a-2)z_{uv} + az_{vv},   \\
            \frac{\partial^2 z}{\partial y^2}          & = \left( -2\frac{\partial}{\partial u} + a\frac{\partial}{\partial v} \right) \left( -2z_u + az_v \right) = 4z_{uu} - 4az_{uv} + a^2 z_{vv}.
        \end{aligned}
    \end{equation}
    代入原方程 $6z_{xx} + z_{xy} - z_{yy} = 0$ 得：
    \begin{equation}
        6(z_{uu} + 2z_{uv} + z_{vv}) + (-2z_{uu} + (a-2)z_{uv} + az_{vv}) - (4z_{uu} - 4az_{uv} + a^2 z_{vv}) = 0.
    \end{equation}
    整理合并同类项：
    \begin{equation}
        (6 - 2 - 4)z_{uu} + (12 + a - 2 + 4a)z_{uv} + (6 + a - a^2)z_{vv} = 0.
    \end{equation}
    即：
    \begin{equation}
        (5a + 10)z_{uv} + (6 + a - a^2)z_{vv} = 0.
    \end{equation}
    要使方程简化为 $z_{uv} = 0$，需满足系数 $6 + a - a^2 = 0$ 且 $5a + 10 \neq 0$．解得：
    \begin{equation}
        a^2 - a - 6 = (a - 3)(a + 2) = 0 \implies a = 3 \text{ 或 } a = -2.
    \end{equation}
    由于当 $a = -2$ 时，变换的雅可比行列式为零且 $5a + 10 = 0$ 导致方程失效，故舍去．取 $a = 3$．
    \begin{equation}
        \boxed{a = 3}
    \end{equation}
\end{solution}


\begin{problem}{偏微分方程}{First-Order-PDE}
    求方程 $\frac{\partial z}{\partial y} = x^2 + 2y$ 满足条件 $z(x, x^2) = 1$ 的解 $z = z(x, y)$．
\end{problem}

\begin{solution}
    对已知方程关于 $y$ 进行积分，得
    \begin{equation}
        z(x, y) = \int (x^2 + 2y) \d y = x^2 y + y^2 + \varphi(x)
    \end{equation}
    其中 $\varphi(x)$ 为待定函数．代入边界条件 $z(x, x^2) = 1$，有
    \begin{equation}
        x^2 \cdot x^2 + (x^2)^2 + \varphi(x) = 1 \implies 2x^4 + \varphi(x) = 1
    \end{equation}
    解得 $\varphi(x) = 1 - 2x^4$．将其代入通解表达式，得到
    \begin{equation}
        \boxed{z(x, y) = x^2 y + y^2 + 1 - 2x^4}
    \end{equation}
\end{solution}


\begin{problem}{复合函数求导}{Partial-Derivative-Chain-Rule-2}
    设 $u = u(x, y)$，当 $y = x^2$ 时有 $u = 1, \frac{\partial u}{\partial x} = x$，求当 $y = x^2$ 时的 $\frac{\partial u}{\partial y}$．
\end{problem}

\begin{solution}
    已知当 $y = x^2$ 时，$u(x, x^2) = 1$．方程两端对 $x$ 求全导数，应用复合函数链式法则得
    \begin{equation}
        \frac{\d}{\d x} [u(x, x^2)] = \frac{\partial u}{\partial x} \cdot \frac{\d x}{\d x} + \frac{\partial u}{\partial y} \cdot \frac{\d (x^2)}{\d x} = 0
    \end{equation}
    整理得
    \begin{equation}
        \frac{\partial u}{\partial x} + 2x \frac{\partial u}{\partial y} = 0
    \end{equation}
    将已知条件 $\frac{\partial u}{\partial x} = x$ 代入，得
    \begin{equation}
        x + 2x \frac{\partial u}{\partial y} = 0
    \end{equation}
    当 $x \neq 0$ 时，解得
    \begin{equation}
        \boxed{\frac{\partial u}{\partial y} = -\frac{1}{2}}
    \end{equation}
\end{solution}


\begin{problem}{二阶偏导数计算}{Second-Order-Partial-Derivatives}
    设 $u = u(x, y)$ 满足方程 $\frac{\partial^2 u}{\partial x^2} - \frac{\partial^2 u}{\partial y^2} = 0$ 以及条件 $u(x, 2x) = x, u'_x(x, 2x) = x^2$，求 $u''_{xx}(x, 2x), u''_{xy}(x, 2x), u''_{yy}(x, 2x)$．（其中二阶偏导数均连续）
\end{problem}

\begin{solution}
    由已知 PDE 知，在任意点处均满足 $u''_{xx} = u''_{yy}$，故在 $(x, 2x)$ 处亦有
    \begin{equation}\label{eq:pde-1}
        u''_{xx}(x, 2x) = u''_{yy}(x, 2x)
    \end{equation}
    对 $u(x, 2x) = x$ 关于 $x$ 求导，得
    \begin{equation}
        u'_x(x, 2x) + 2u'_y(x, 2x) = 1
    \end{equation}
    将 $u'_x(x, 2x) = x^2$ 代入上式，得 $2u'_y(x, 2x) = 1 - x^2$，即
    \begin{equation}\label{eq:uy-val}
        u'_y(x, 2x) = \frac{1 - x^2}{2}
    \end{equation}
    将 $u'_x(x, 2x) = x^2$ 两端关于 $x$ 求导，利用链式法则有
    \begin{equation}\label{eq:uxx-1}
        u''_{xx}(x, 2x) + 2u''_{xy}(x, 2x) = 2x
    \end{equation}
    对 \zcref{eq:uy-val} 两端关于 $x$ 求导，得
    \begin{equation}
        u''_{yx}(x, 2x) + 2u''_{yy}(x, 2x) = -x
    \end{equation}
    由于二阶偏导数连续，则 $u''_{xy} = u''_{yx}$，结合 \zcref{eq:pde-1} 得
    \begin{equation}\label{eq:uxx-2}
        u''_{xy}(x, 2x) + 2u''_{xx}(x, 2x) = -x
    \end{equation}
    联立 \zcref{eq:uxx-1} 与 \zcref{eq:uxx-2} 构建方程组：
    \begin{equation}
        \begin{cases}
            u''_{xx} + 2u''_{xy} = 2x \\
            2u''_{xx} + u''_{xy} = -x
        \end{cases}
    \end{equation}
    解得 $u''_{xx} = -\frac{4}{3}x$，$u''_{xy} = \frac{5}{3}x$．由 \zcref{eq:pde-1} 可得 $u''_{yy} = -\frac{4}{3}x$．
    \begin{equation}
        \boxed{u''_{xx}(x, 2x) = -\frac{4}{3}x, \quad u''_{xy}(x, 2x) = \frac{5}{3}x, \quad u''_{yy}(x, 2x) = -\frac{4}{3}x}
    \end{equation}
\end{solution}


\begin{problem}{复合函数微分}{Differential-Composite-Functions}
    求下列复合函数的微分 $\d u$：
    \begin{enumerate}
        \item $u = f(t)$, $t = x + y$；
        \item $u = f(\xi, \eta)$, $\xi = xy$, $\eta = \frac{x}{y}$；
        \item $u = f(x, y, z)$, $x = t$, $y = t^2$, $z = t^3$；
        \item $u = f(x, \xi, \eta)$, $\xi = x^2 + y^2$, $\eta = x^2 + y^2 + z^2$；
        \item $u = f(\xi, \eta, \zeta)$, $\xi = x^2 + y^2$, $\eta = x^2 - y^2$, $\zeta = 2xy$.
    \end{enumerate}
\end{problem}

\begin{solution}
    \ansitem{1}

    \begin{equation}
        \d u = f'(t) \d t = f'(t) (\d x + \d y).
    \end{equation}
    \begin{equation}
        \boxed{\d u = f'(t) (\d x + \d y)}
    \end{equation}

    \ansitem{2}

    \begin{equation}
        \begin{aligned}
            \d u & = f_{\xi} \d \xi + f_{\eta} \d \eta                                                                              \\
                 & = f_{\xi} (y \d x + x \d y) + f_{\eta} \frac{y \d x - x \d y}{y^2}                                               \\
                 & = \left( y f_{\xi} + \frac{1}{y} f_{\eta} \right) \d x + \left( x f_{\xi} - \frac{x}{y^2} f_{\eta} \right) \d y.
        \end{aligned}
    \end{equation}
    \begin{equation}
        \boxed{\d u = \left( y f_{\xi} + \frac{1}{y} f_{\eta} \right) \d x + \left( x f_{\xi} - \frac{x}{y^2} f_{\eta} \right) \d y}
    \end{equation}

    \ansitem{3}

    \begin{equation}
        \begin{aligned}
            \d u & = f_x \d x + f_y \d y + f_z \d z             \\
                 & = f_x \d t + f_y (2t \d t) + f_z (3t^2 \d t) \\
                 & = (f_x + 2t f_y + 3t^2 f_z) \d t.
        \end{aligned}
    \end{equation}
    \begin{equation}
        \boxed{\d u = (f_x + 2t f_y + 3t^2 f_z) \d t}
    \end{equation}

    \ansitem{4}

    \begin{equation}
        \begin{aligned}
            \d u & = f_1 \d x + f_2 \d \xi + f_3 \d \eta                                    \\
                 & = f_1 \d x + f_2 (2x \d x + 2y \d y) + f_3 (2x \d x + 2y \d y + 2z \d z) \\
                 & = (f_1 + 2x f_2 + 2x f_3) \d x + (2y f_2 + 2y f_3) \d y + 2z f_3 \d z.
        \end{aligned}
    \end{equation}
    \begin{equation}
        \boxed{\d u = (f_x + 2x f_{\xi} + 2x f_{\eta}) \d x + (2y f_{\xi} + 2y f_{\eta}) \d y + 2z f_{\eta} \d z}
    \end{equation}

    \ansitem{5}

    \begin{equation}
        \begin{aligned}
            \d u & = f_{\xi} \d \xi + f_{\eta} \d \eta + f_{\zeta} \d \zeta                                           \\
                 & = f_{\xi} (2x \d x + 2y \d y) + f_{\eta} (2x \d x - 2y \d y) + f_{\zeta} (2y \d x + 2x \d y)       \\
                 & = (2x f_{\xi} + 2x f_{\eta} + 2y f_{\zeta}) \d x + (2y f_{\xi} - 2y f_{\eta} + 2x f_{\zeta}) \d y.
        \end{aligned}
    \end{equation}
    \begin{equation}
        \boxed{\d u = 2(x f_{\xi} + x f_{\eta} + y f_{\zeta}) \d x + 2(y f_{\xi} - y f_{\eta} + x f_{\zeta}) \d y}
    \end{equation}
\end{solution}

\begin{problem}{球坐标下拉普拉斯方程}{Spherical-Laplace-Equation}
    求在球坐标下的 Laplace 方程的形式．
\end{problem}

\begin{solution}
    在球坐标系 $(r, \theta, \varphi)$ 下，其中 $r$ 为径向距离，$\theta$ 为极角（余纬角），$\varphi$ 为方位角．坐标变换为：
    \begin{equation}
        \left\{
        \begin{aligned}
            x & = r \sin \theta \cos \varphi, \\
            y & = r \sin \theta \sin \varphi, \\
            z & = r \cos \theta.
        \end{aligned}
        \right.
    \end{equation}
    利用链式法则进行坐标变换推导，Laplace 算子 $\nabla^2 u = \frac{\partial^2 u}{\partial x^2} + \frac{\partial^2 u}{\partial y^2} + \frac{\partial^2 u}{\partial z^2}$ 在该坐标系下的形式为：
    \begin{equation}
        \nabla^2 u = \frac{1}{r^2} \frac{\partial}{\partial r} \left( r^2 \frac{\partial u}{\partial r} \right) + \frac{1}{r^2 \sin \theta} \frac{\partial}{\partial \theta} \left( \sin \theta \frac{\partial u}{\partial \theta} \right) + \frac{1}{r^2 \sin^2 \theta} \frac{\partial^2 u}{\partial \varphi^2} = 0.
    \end{equation}
    \begin{equation}
        \boxed{\frac{1}{r^2} \frac{\partial}{\partial r} \left( r^2 \frac{\partial u}{\partial r} \right) + \frac{1}{r^2 \sin \theta} \frac{\partial}{\partial \theta} \left( \sin \theta \frac{\partial u}{\partial \theta} \right) + \frac{1}{r^2 \sin^2 \theta} \frac{\partial^2 u}{\partial \varphi^2} = 0}
    \end{equation}
\end{solution}

\begin{problem}{雅可比行列式}{Jacobian-Determinant}
    求直角坐标和极坐标的坐标变换 $x = x(r, \theta) = r \cos \theta, y = y(r, \theta) = r \sin \theta$ 的 Jacobi 行列式．
\end{problem}

\begin{solution}
    根据 Jacobi 行列式的定义，变换的行列式 $J$ 为：
    \begin{equation}
        J = \frac{\partial(x, y)}{\partial(r, \theta)} = \begin{vmatrix}
            \frac{\partial x}{\partial r} & \frac{\partial x}{\partial \theta} \\
            \frac{\partial y}{\partial r} & \frac{\partial y}{\partial \theta}
        \end{vmatrix}.
    \end{equation}
    计算各偏导数：
    \begin{equation}
        \frac{\partial x}{\partial r} = \cos \theta, \quad \frac{\partial x}{\partial \theta} = -r \sin \theta, \quad \frac{\partial y}{\partial r} = \sin \theta, \quad \frac{\partial y}{\partial \theta} = r \cos \theta.
    \end{equation}
    代入行列式得：
    \begin{equation}
        J = \begin{vmatrix}
            \cos \theta & -r \sin \theta \\
            \sin \theta & r \cos \theta
        \end{vmatrix} = r \cos^2 \theta - (-r \sin^2 \theta) = r(\cos^2 \theta + \sin^2 \theta) = r.
    \end{equation}
    \begin{equation}
        \boxed{J = r}
    \end{equation}
\end{solution}

\endinput
